\chapter{Stack and Queue}
\chapterepigraph{A stack remembers what just happened. A queue remembers
what happened first. Almost every problem in this chapter is, underneath,
a decision about which of these two memories the problem actually needs --
and a surprising number of ``queue'' problems turn out to be built from
stacks, and vice versa.}

\section{Stacks, Queues, and the Monotonic Stack Idea}
\label{sec:sq-intro}

A \textbf{stack} supports two operations, \textit{push} and \textit{pop},
with Last-In-First-Out (LIFO) order: the most recently pushed element is
always the first one popped. A \textbf{queue} supports \textit{enqueue}
and \textit{dequeue} with First-In-First-Out (FIFO) order: the earliest
enqueued element still present is always the first one dequeued. Both can
be implemented on top of an array, a linked list, or -- as the first
section of this chapter shows -- on top of \emph{each other}.

\begin{patternbox}[How to Recognise Which One You Need]
\begin{itemize}
  \item \textbf{``Matching/nesting/balanced''} (parentheses, expression
        evaluation, undo history) -- a stack, because the most recent
        unmatched thing is always the first thing that needs to be closed
        or undone.
  \item \textbf{``Next/previous greater/smaller element''} -- a
        \textbf{monotonic stack}: a stack kept strictly increasing or
        decreasing from bottom to top by popping elements that violate
        the order before pushing a new one. This is the single most
        important pattern in this chapter, covered in depth starting at
        Section~\ref{sec:sq-next-greater-1}.
  \item \textbf{``Maximum/minimum over every sliding window''} -- a
        \textbf{monotonic deque} (Section~\ref{sec:sq-sliding-window-max}):
        a double-ended queue that behaves like a monotonic stack at one
        end, while also supporting eviction from the other end as the
        window slides.
  \item \textbf{``Most/least recently used,'' with fast lookup} -- a
        combination of a hash map (for $O(1)$ lookup) and a doubly linked
        list or queue-like ordering (for fast reordering by recency), as
        in LRU/LFU Cache (Sections~\ref{sec:sq-lru-cache}--\ref{sec:sq-lfu-cache}).
\end{itemize}
\end{patternbox}

\subsection*{The Monotonic Stack, Explained Once}

\begin{ideabox}[Why a Monotonic Stack Finds ``Next Greater'' in $O(n)$]
To find, for every element, the next element to its right that is
strictly greater, scan from \textbf{right to left}, maintaining a stack
that is kept strictly \emph{decreasing} from bottom to top. Before
processing element $i$, pop every stack element that is $\le
\textit{arr}[i]$ -- those elements can never again be the ``next greater''
answer for anything to their left, because $\textit{arr}[i]$ is closer and
at least as large, so it would always be found first; equivalently they have
just been proven useless once $\textit{arr}[i]$ shows up. After popping,
whatever remains on top of the stack (if anything) is the next greater
element for $\textit{arr}[i]$ -- the nearest still-undefeated candidate
to $i$'s right. Then push $\textit{arr}[i]$ itself, since it might be the
answer for elements further to the left.

The reason this is $O(n)$ total, not $O(n^2)$: every element is pushed
exactly once and popped at most once, across the \emph{entire} scan -- so
even though some steps pop many elements, the total number of pops across
all steps combined can never exceed $n$.
\end{ideabox}

This single idea -- maintain a monotonic stack, popping whatever the
current element invalidates before pushing -- powers roughly a third of
this chapter's problems, sometimes directly (next/previous greater or
smaller element) and sometimes as a hidden sub-step inside a larger
problem (histogram area, trapping rain water, sum of subarray minimums).

Each section in this chapter follows the same structure as the rest of
this book: problem statement and keywords; the reasoning that identifies
the right stack/queue technique and why it is correct; the algorithm in
words and pseudocode; a hand-traced dry run; complete compilable C++ code;
and a time/space complexity analysis.

\newpage

\section{Implement a Stack using an Array}
\label{sec:sq-stack-array}

\problemstatement{Implement a stack supporting $\textit{push}(x)$,
$\textit{pop}()$, $\textit{top}()$, and $\textit{isEmpty}()$, backed by a
fixed-size array, each operation running in $O(1)$.}

\begin{patternbox}
No pattern recognition is needed here -- this is foundational machinery.
The only idea worth internalizing is that a stack needs just \textbf{one}
extra piece of state beyond the array itself: an index, conventionally
called \textit{top}, pointing at the most recently pushed element.
\end{patternbox}

\subsection*{Understanding the problem carefully}

Since a stack only ever adds or removes from one end (the ``top''), an
array can simulate it perfectly by treating one end (conventionally the
highest used index) as that top, and tracking that index explicitly. No
shifting of other elements is ever required, which is exactly what keeps
every operation $O(1)$.

\begin{ideabox}[The Single Pointer Invariant]
Maintain $\textit{top}$ as the index of the most recently pushed element
($-1$ when empty). $\textit{push}(x)$ writes $x$ at index $\textit{top}+1$
and increments \textit{top}. $\textit{pop}()$ reads $\textit{arr}
[\textit{top}]$ and decrements \textit{top} (the old value is left in
place but is considered logically removed, since future pushes will
overwrite it).
\end{ideabox}

\subsection*{Algorithm}

\begin{enumerate}
  \item \textit{top} $\gets -1$ initially.
  \item $\textit{push}(x)$: increment \textit{top}; set
        $\textit{arr}[\textit{top}] \gets x$.
  \item $\textit{pop}()$: read $\textit{arr}[\textit{top}]$; decrement
        \textit{top}; return the read value.
  \item $\textit{top}()$: return $\textit{arr}[\textit{top}]$.
  \item $\textit{isEmpty}()$: return $(\textit{top} = -1)$.
\end{enumerate}

\subsection*{Dry run}

\dryrun{Push $10, 20, 30$; pop once; push $40$.}

\begin{itemize}
  \item Push $10$: \textit{top}$=0$, $\textit{arr}=[10]$.
  \item Push $20$: \textit{top}$=1$, $\textit{arr}=[10,20]$.
  \item Push $30$: \textit{top}$=2$, $\textit{arr}=[10,20,30]$.
  \item Pop: returns $30$; \textit{top}$=1$.
  \item Push $40$: \textit{top}$=2$, $\textit{arr}=[10,20,40]$ (index $2$
        is overwritten).
\end{itemize}

\subsection*{C++ implementation}

\begin{lstlisting}
#include <bits/stdc++.h>
using namespace std;

class Stack {
    vector<int> arr;
    int topIdx;

public:
    Stack(int capacity) : arr(capacity), topIdx(-1) {}

    void push(int x) {
        arr[++topIdx] = x;
    }

    int pop() {
        return arr[topIdx--];
    }

    int top() {
        return arr[topIdx];
    }

    bool isEmpty() {
        return topIdx == -1;
    }
};

int main() {
    Stack s(10);
    s.push(10); s.push(20); s.push(30);
    cout << "Popped: " << s.pop() << endl;
    s.push(40);
    cout << "Top: " << s.top() << endl;
    return 0;
}
\end{lstlisting}

\begin{complexitybox}
\textbf{Time Complexity.} $O(1)$ for every operation -- only a single
index update and a single array access are involved.
\textbf{Space Complexity.} $O(\textit{capacity})$ for the underlying
array.
\end{complexitybox}

\begin{takeawaybox}
A stack needs exactly one pointer into a contiguous array. Compare this to
Section~\ref{sec:sq-queue-linked-list}, where a queue's two-ended access
pattern makes a plain array awkward (without wraparound) and motivates a
linked-list-based implementation instead.
\end{takeawaybox}

\section{Implement a Stack using a Linked List}
\label{sec:sq-stack-linked-list}

\problemstatement{Implement a stack supporting $\textit{push}(x)$,
$\textit{pop}()$, $\textit{top}()$, and $\textit{isEmpty}()$, backed by a
singly linked list, each operation running in $O(1)$, with no fixed
capacity limit.}

\begin{patternbox}
The keyword shift from Section~\ref{sec:sq-stack-array} is \emph{``no
fixed capacity''} -- an array-backed stack needs to know its maximum size
up front (or pay for resizing), while a linked list grows and shrinks node
by node with no such limit. The structural insight: a stack's LIFO order
matches a linked list's \textbf{head} perfectly, since inserting and
removing at the head of a singly linked list are both already $O(1)$.
\end{patternbox}

\subsection*{Understanding the problem carefully}

If we always insert new elements at the \textbf{head} of the list and
always remove from the head, the most recently inserted element is always
the first one removed -- exactly LIFO order, with no need to traverse the
list at all.

\begin{ideabox}[Head Insertion = Push, Head Removal = Pop]
Maintain a single pointer \textit{head} to the front of the list (null
when empty). $\textit{push}(x)$ creates a new node with $\textit{next}
\gets \textit{head}$, then sets $\textit{head}$ to this new node.
$\textit{pop}()$ reads $\textit{head}$'s value, then advances
$\textit{head} \gets \textit{head}.\textit{next}$.
\end{ideabox}

\subsection*{Algorithm}

\begin{enumerate}
  \item $\textit{head} \gets$ null initially.
  \item $\textit{push}(x)$: $\textit{node} \gets$ new node with value $x$,
        $\textit{next} \gets \textit{head}$; $\textit{head} \gets
        \textit{node}$.
  \item $\textit{pop}()$: $\textit{val} \gets \textit{head}.\textit{val}$;
        $\textit{head} \gets \textit{head}.\textit{next}$; return
        \textit{val}.
  \item $\textit{top}()$: return $\textit{head}.\textit{val}$.
  \item $\textit{isEmpty}()$: return $(\textit{head} = \text{null})$.
\end{enumerate}

\subsection*{Dry run}

\dryrun{Push $1,2,3$; pop once.}

\begin{itemize}
  \item Push $1$: list $= 1 \to \text{null}$.
  \item Push $2$: list $= 2 \to 1 \to \text{null}$.
  \item Push $3$: list $= 3 \to 2 \to 1 \to \text{null}$.
  \item Pop: returns $3$; list $= 2 \to 1 \to \text{null}$.
\end{itemize}

\subsection*{C++ implementation}

\begin{lstlisting}
#include <bits/stdc++.h>
using namespace std;

class Stack {
    struct Node {
        int val;
        Node* next;
        Node(int v, Node* n) : val(v), next(n) {}
    };
    Node* head = nullptr;

public:
    void push(int x) {
        head = new Node(x, head);
    }

    int pop() {
        int val = head->val;
        Node* old = head;
        head = head->next;
        delete old;
        return val;
    }

    int top() {
        return head->val;
    }

    bool isEmpty() {
        return head == nullptr;
    }
};

int main() {
    Stack s;
    s.push(1); s.push(2); s.push(3);
    cout << "Popped: " << s.pop() << endl;
    cout << "Top: " << s.top() << endl;
    return 0;
}
\end{lstlisting}

\begin{complexitybox}
\textbf{Time Complexity.} $O(1)$ for every operation -- all of them touch
only the head node.
\textbf{Space Complexity.} $O(n)$ for $n$ stored elements, with one node
of constant-size overhead each, and no wasted unused capacity (unlike a
fixed-size array implementation).
\end{complexitybox}

\begin{takeawaybox}
A stack's LIFO discipline matches a singly linked list's head exactly --
no \textit{tail} pointer is ever needed, since insertion and removal both
happen at the same end. This is the cleanest possible linked-list data
structure to implement, and it is exactly why the next section's queue
(needing access at \emph{both} ends) is structurally harder.
\end{takeawaybox}

\section{Implement a Queue using a Linked List}
\label{sec:sq-queue-linked-list}

\problemstatement{Implement a queue supporting $\textit{enqueue}(x)$,
$\textit{dequeue}()$, $\textit{front}()$, and $\textit{isEmpty}()$, backed
by a singly linked list, each operation running in $O(1)$.}

\begin{patternbox}
A queue's FIFO discipline needs insertion at \textbf{one} end and removal
from the \textbf{other} -- unlike a stack, a single \textit{head} pointer
is not enough, because removing from the head and inserting at the head
would give LIFO, not FIFO. The keyword to notice: this is the first
structure in this chapter that genuinely requires \textbf{two} pointers
into the list (\textit{head} and \textit{tail}) to stay $O(1)$.
\end{patternbox}

\subsection*{Understanding the problem carefully}

If we insert new elements at the \textbf{tail} and remove from the
\textbf{head}, the earliest-inserted, not-yet-removed element is always
at the head, ready to be dequeued -- exactly FIFO order. Maintaining only
a \textit{head} pointer would make tail-insertion cost $O(n)$ (we would
have to walk the whole list to find the current last node), so we keep an
explicit \textit{tail} pointer too, updated on every enqueue.

\begin{ideabox}[Two Pointers: Head for Removal, Tail for Insertion]
$\textit{enqueue}(x)$ creates a new node and links it after the current
\textit{tail} (or sets it as both \textit{head} and \textit{tail} if the
queue was empty), then updates \textit{tail} to point to this new node.
$\textit{dequeue}()$ reads \textit{head}'s value, then advances
$\textit{head} \gets \textit{head}.\textit{next}$ (and resets
\textit{tail} to null if the queue becomes empty).
\end{ideabox}

\subsection*{Algorithm}

\begin{enumerate}
  \item $\textit{head}, \textit{tail} \gets$ null initially.
  \item $\textit{enqueue}(x)$: $\textit{node} \gets$ new node with value
        $x$. If $\textit{tail} = $ null: $\textit{head} \gets
        \textit{tail} \gets \textit{node}$. Else:
        $\textit{tail}.\textit{next} \gets \textit{node}$;
        $\textit{tail} \gets \textit{node}$.
  \item $\textit{dequeue}()$: $\textit{val} \gets
        \textit{head}.\textit{val}$; $\textit{head} \gets
        \textit{head}.\textit{next}$; if $\textit{head} = $ null:
        $\textit{tail} \gets$ null; return \textit{val}.
  \item $\textit{front}()$: return $\textit{head}.\textit{val}$.
  \item $\textit{isEmpty}()$: return $(\textit{head} = \text{null})$.
\end{enumerate}

\subsection*{Dry run}

\dryrun{Enqueue $1,2,3$; dequeue once.}

\begin{itemize}
  \item Enqueue $1$: $\textit{head}=\textit{tail}=1$.
  \item Enqueue $2$: $1 \to 2$; $\textit{head}=1$, $\textit{tail}=2$.
  \item Enqueue $3$: $1\to2\to3$; $\textit{head}=1$, $\textit{tail}=3$.
  \item Dequeue: returns $1$; $\textit{head}=2$ now; list $2\to3$.
\end{itemize}

\subsection*{C++ implementation}

\begin{lstlisting}
#include <bits/stdc++.h>
using namespace std;

class Queue {
    struct Node {
        int val;
        Node* next;
        Node(int v) : val(v), next(nullptr) {}
    };
    Node* head = nullptr;
    Node* tail = nullptr;

public:
    void enqueue(int x) {
        Node* node = new Node(x);
        if (!tail) {
            head = tail = node;
        } else {
            tail->next = node;
            tail = node;
        }
    }

    int dequeue() {
        int val = head->val;
        Node* old = head;
        head = head->next;
        if (!head) tail = nullptr;
        delete old;
        return val;
    }

    int front() {
        return head->val;
    }

    bool isEmpty() {
        return head == nullptr;
    }
};

int main() {
    Queue q;
    q.enqueue(1); q.enqueue(2); q.enqueue(3);
    cout << "Dequeued: " << q.dequeue() << endl;
    cout << "Front: " << q.front() << endl;
    return 0;
}
\end{lstlisting}

\begin{complexitybox}
\textbf{Time Complexity.} $O(1)$ for every operation, since \textit{tail}
removes the need to traverse the list on enqueue.
\textbf{Space Complexity.} $O(n)$ for $n$ stored elements.
\end{complexitybox}

\begin{takeawaybox}
The need for two pointers (\textit{head} and \textit{tail}) rather than
one is the structural signature of FIFO order on a singly linked list.
This same head/tail pairing reappears, in a different form, inside the
doubly linked list used by LRU Cache (Section~\ref{sec:sq-lru-cache}).
\end{takeawaybox}

\section{Implement a Stack using a Queue}
\label{sec:sq-stack-using-queue}

\problemstatement{Implement a stack (LIFO: \textit{push}, \textit{pop},
\textit{top}) using only the standard operations of a queue
(\textit{enqueue}, \textit{dequeue}, \textit{front}) as building blocks --
no direct array or linked-list indexing allowed.}

\begin{patternbox}
\emph{``Implement structure $X$ using only structure $Y$'s operations''}
is a request to simulate one access order using a tool built for the
opposite access order. The trick, here and in
Section~\ref{sec:sq-queue-using-stack}, is realizing that you can pay an
\textbf{extra cost at one specific operation} (either push or pop, your
choice) to physically reorder the underlying structure's contents,
restoring the access pattern you actually need.
\end{patternbox}

\subsection*{Understanding the problem carefully}

A queue always exposes its \emph{oldest} element at the front. A stack
needs to expose its \emph{newest} element at the top. If, every time we
push a new element, we immediately rotate the entire queue so that the
new element ends up at the front (instead of the back, where a plain
enqueue would place it), then the front of the queue always holds the most
recently pushed element -- exactly what a stack's \textit{top} needs.

\begin{ideabox}[Make Push Expensive, Keep Pop/Top Cheap]
$\textit{push}(x)$: enqueue $x$ normally (it lands at the back). Then
rotate the queue by dequeuing and re-enqueuing every \emph{other} element
(everything that was already in the queue before $x$ was added) -- this
moves $x$, one rotation step at a time, all the way around to the front.
$\textit{pop}()$ and $\textit{top}()$ then become trivial: simply dequeue
(or peek) the queue's front, which is always the most recently pushed
element by the invariant just established.
\end{ideabox}

\subsection*{Why rotating after every push maintains the LIFO invariant}

Each push performs exactly $(\textit{size before this push})$ rotation
steps, which is enough to move the newly enqueued element from the back
to the front while preserving the relative order of everything else
(each older element is dequeued and immediately re-enqueued once, so
their relative order among themselves never changes). After this
rotation, the queue's front-to-back order is exactly ``most recently
pushed first,'' i.e.\ LIFO order, by induction: if this invariant held
before the push, the rotation restores it after the push as well.

\subsection*{Algorithm}

\begin{enumerate}
  \item $\textit{push}(x)$: let $s$ be the queue's current size before
        this push. Enqueue $x$. Repeat $s$ times: dequeue an element and
        immediately re-enqueue it.
  \item $\textit{pop}()$: dequeue and return the front.
  \item $\textit{top}()$: return the front, without removing it.
\end{enumerate}

\subsection*{Dry run}

\dryrun{Push $1, 2, 3$ in order.}

\begin{itemize}
  \item Push $1$: queue $=[1]$ ($s=0$, no rotation needed).
  \item Push $2$: enqueue $2$: queue $=[1,2]$. Rotate $s=1$ time:
        dequeue $1$, enqueue $1$: queue $=[2,1]$.
  \item Push $3$: enqueue $3$: queue $=[2,1,3]$. Rotate $s=2$ times:
        dequeue $2$, enqueue $2$: queue $=[1,3,2]$; dequeue $1$, enqueue
        $1$: queue $=[3,2,1]$.
\end{itemize}

The queue's front-to-back order is $[3,2,1]$ -- exactly LIFO order, most
recently pushed ($3$) at the front.

\subsection*{C++ implementation}

\begin{lstlisting}
#include <bits/stdc++.h>
using namespace std;

class StackUsingQueue {
    queue<int> q;

public:
    void push(int x) {
        int s = q.size();
        q.push(x);
        for (int i = 0; i < s; i++) {
            q.push(q.front());
            q.pop();
        }
    }

    int pop() {
        int val = q.front();
        q.pop();
        return val;
    }

    int top() {
        return q.front();
    }
};

int main() {
    StackUsingQueue s;
    s.push(1); s.push(2); s.push(3);
    cout << "Top: " << s.top() << endl;
    cout << "Popped: " << s.pop() << endl;
    return 0;
}
\end{lstlisting}

\begin{complexitybox}
\textbf{Time Complexity.} $\textit{push}$ costs $O(n)$ (the rotation
touches every existing element). $\textit{pop}$ and $\textit{top}$ cost
$O(1)$.
\textbf{Space Complexity.} $O(n)$ for the underlying queue.
\end{complexitybox}

\begin{takeawaybox}
When forced to simulate one access order using a structure built for the
opposite order, you must pay an extra reordering cost \emph{somewhere} --
the only design choice is \emph{which} operation absorbs that cost. Here
we chose to make \textit{push} expensive and \textit{pop}/\textit{top}
cheap; Section~\ref{sec:sq-queue-using-stack} makes the opposite choice
for the reverse simulation, illustrating both options.
\end{takeawaybox}

\section{Implement a Queue using Stacks}
\label{sec:sq-queue-using-stack}

\problemstatement{Implement a queue (FIFO: \textit{enqueue},
\textit{dequeue}, \textit{front}) using only the standard operations of a
stack (\textit{push}, \textit{pop}, \textit{top}) as building blocks.}

\begin{patternbox}
The reverse simulation of Section~\ref{sec:sq-stack-using-queue}. The key
fact to recall: reversing a sequence \textbf{twice} restores its original
order -- so pouring elements through \emph{two} stacks in succession
(stack $\to$ stack $\to$ stack) flips LIFO order into FIFO order, because
each individual pour through a stack reverses the order once.
\end{patternbox}

\subsection*{Understanding the problem carefully}

A single stack inherently reverses whatever order you feed into it: if
you push $1,2,3$ and then pop everything, you get $3,2,1$ -- reversed. If
we use \textbf{two} stacks, call them \textit{inStack} and
\textit{outStack}, and whenever \textit{outStack} runs empty we pour all
of \textit{inStack} into it (popping from \textit{inStack} and pushing
each onto \textit{outStack}), the act of pouring reverses the order a
\emph{second} time -- and two reversals cancel out, restoring the
original (FIFO) order on \textit{outStack}'s pop sequence.

\begin{ideabox}[Two Stacks: One for Input, One for Output]
$\textit{enqueue}(x)$: simply push $x$ onto \textit{inStack} -- $O(1)$,
no reordering needed yet. $\textit{dequeue}()$ (and $\textit{front}()$):
if \textit{outStack} is empty, pour the entirety of \textit{inStack} into
it (pop each element off \textit{inStack}, push it onto \textit{outStack});
this reverses \textit{inStack}'s order, placing the
\emph{oldest-enqueued, not-yet-dequeued} element on top of
\textit{outStack}. Then pop (or peek) \textit{outStack}'s top.
\end{ideabox}

\subsection*{Why pouring only when \textit{outStack} is empty gives amortized $O(1)$}

If we poured on \emph{every} dequeue, each dequeue would cost $O(n)$.
Instead, we only pour when \textit{outStack} is completely empty -- and
once poured, \textit{outStack} correctly holds a whole batch of elements
in the right FIFO order, which can each be dequeued in $O(1)$ without
needing another pour until that batch is exhausted. Each individual
element is poured (moved from \textit{inStack} to \textit{outStack}) at
most \emph{once} in its entire lifetime, so across any sequence of $n$
enqueue/dequeue operations, the total cost of all pours combined is
$O(n)$, giving $O(1)$ \emph{amortized} time per operation even though some
individual dequeues are more expensive than others.

\subsection*{Algorithm}

\begin{enumerate}
  \item $\textit{enqueue}(x)$: push $x$ onto \textit{inStack}.
  \item $\textit{dequeue}()$ / $\textit{front}()$: if \textit{outStack}
        is empty, repeatedly pop from \textit{inStack} and push onto
        \textit{outStack} until \textit{inStack} is empty. Then pop (or
        peek) \textit{outStack}.
\end{enumerate}

\subsection*{Dry run}

\dryrun{Enqueue $1,2,3$; dequeue twice; enqueue $4$; dequeue once.}

\begin{itemize}
  \item Enqueue $1,2,3$: \textit{inStack} (bottom to top) $=[1,2,3]$;
        \textit{outStack} $=[]$.
  \item Dequeue: \textit{outStack} empty, so pour: pop $3,2,1$ from
        \textit{inStack}, push onto \textit{outStack}, giving
        \textit{outStack} (bottom to top) $=[3,2,1]$. Pop \textit{outStack}'s
        top: returns $1$. \textit{outStack}$=[3,2]$.
  \item Dequeue: \textit{outStack} not empty ($[3,2]$). Pop top: returns
        $2$. \textit{outStack}$=[3]$.
  \item Enqueue $4$: push onto \textit{inStack}$=[4]$.
        \textit{outStack} still $=[3]$, untouched.
  \item Dequeue: \textit{outStack} not empty. Pop top: returns $3$.
\end{itemize}

The dequeue sequence $1,2,3$ correctly matches the enqueue order.

\subsection*{C++ implementation}

\begin{lstlisting}
#include <bits/stdc++.h>
using namespace std;

class QueueUsingStack {
    stack<int> inStack, outStack;

    void transferIfNeeded() {
        if (outStack.empty()) {
            while (!inStack.empty()) {
                outStack.push(inStack.top());
                inStack.pop();
            }
        }
    }

public:
    void enqueue(int x) {
        inStack.push(x);
    }

    int dequeue() {
        transferIfNeeded();
        int val = outStack.top();
        outStack.pop();
        return val;
    }

    int front() {
        transferIfNeeded();
        return outStack.top();
    }
};

int main() {
    QueueUsingStack q;
    q.enqueue(1); q.enqueue(2); q.enqueue(3);
    cout << q.dequeue() << " " << q.dequeue() << endl;
    q.enqueue(4);
    cout << q.dequeue() << endl;
    return 0;
}
\end{lstlisting}

\begin{complexitybox}
\textbf{Time Complexity.} $\textit{enqueue}$ is always $O(1)$.
$\textit{dequeue}$ and $\textit{front}$ are $O(1)$ \textbf{amortized}
(occasionally $O(n)$ for a pour, but each element is poured at most once
across its lifetime).
\textbf{Space Complexity.} $O(n)$ across the two stacks combined.
\end{complexitybox}

\begin{takeawaybox}
Reversing twice restores original order -- this is the entire idea behind
using two stacks to simulate a queue, and it is worth remembering as a
general trick: any time an algorithm needs to ``undo'' a reversal it was
forced into, pouring through a second stack is often the simplest fix.
\end{takeawaybox}

\section{Min Stack}
\label{sec:sq-min-stack}

\problemstatement{Design a stack supporting $\textit{push}(x)$,
$\textit{pop}()$, $\textit{top}()$, and $\textit{getMin}()$ -- the last
returning the minimum element currently in the stack -- all in $O(1)$
time.}

\begin{patternbox}
\emph{``Maintain a running aggregate (min, max, sum) that must stay
correct as elements are pushed \emph{and} popped''} is the keyword
pattern here. A plain running variable breaks the moment a \textit{pop}
removes the element that was responsible for the current minimum -- there
would be no way to recover the \emph{previous} minimum without
re-scanning. The fix is to keep a full \textbf{history} of the minimum, one
entry per stack position, so that popping the main stack and popping the
history together always keeps both in sync.
\end{patternbox}

\subsection*{Understanding the problem carefully}

The difficulty is specifically about \textit{pop}: tracking the minimum
as elements are pushed is trivial (compare the new element to the current
minimum). The problem is that after a \textit{pop}, the minimum might need
to \emph{revert} to whatever it was before the popped element was pushed
-- and that historical value must be available instantly, not
recomputed.

\begin{ideabox}[A Second Stack Tracks the Minimum-So-Far at Every Depth]
Maintain an auxiliary \textit{minStack}, parallel to the main stack. On
$\textit{push}(x)$: push $x$ onto the main stack as usual, and push
$\min(x, \textit{minStack.top}())$ onto \textit{minStack} (or just $x$ if
\textit{minStack} is currently empty). On $\textit{pop}()$: pop both
stacks together. $\textit{getMin}()$ simply reads \textit{minStack}'s
top.
\end{ideabox}

\subsection*{Why popping both stacks together keeps the invariant correct}

\textit{minStack}'s $i$-th entry (counting from the bottom) is, by
construction, the minimum of the main stack's bottom $i$ elements at the
moment the $i$-th element was pushed -- and since both stacks always have
the same height and are pushed to and popped from in lockstep,
\textit{minStack}'s current top is always exactly the minimum of
\emph{whatever the main stack currently contains}, not some stale value
from before a pop. This is the same ``parallel history'' idea that makes
undo/redo systems work: rather than computing an aggregate forward only,
store enough history to reconstruct any previous state instantly.

\subsection*{Algorithm}

\begin{enumerate}
  \item $\textit{push}(x)$: push $x$ onto \textit{mainStack}. Push
        $\min(x, \text{current } \textit{getMin}())$ onto
        \textit{minStack} (using $x$ itself if \textit{minStack} was
        empty).
  \item $\textit{pop}()$: pop both \textit{mainStack} and
        \textit{minStack}.
  \item $\textit{top}()$: return \textit{mainStack}'s top.
  \item $\textit{getMin}()$: return \textit{minStack}'s top.
\end{enumerate}

\subsection*{Dry run}

\dryrun{Push $5, 3, 7, 2$; pop once; query getMin.}

\begin{itemize}
  \item Push $5$: main$=[5]$, min$=[5]$.
  \item Push $3$: main$=[5,3]$, min$=[5,\min(3,5){=}3]$.
  \item Push $7$: main$=[5,3,7]$, min$=[5,3,\min(7,3){=}3]$.
  \item Push $2$: main$=[5,3,7,2]$, min$=[5,3,3,\min(2,3){=}2]$.
  \item Pop: both stacks lose their top. main$=[5,3,7]$, min$=[5,3,3]$.
  \item getMin(): top of min stack is $3$.
\end{itemize}

Correct: after removing $2$, the remaining stack $[5,3,7]$ has minimum
$3$.

\subsection*{C++ implementation}

\begin{lstlisting}
#include <bits/stdc++.h>
using namespace std;

class MinStack {
    stack<int> mainStack, minStack;

public:
    void push(int x) {
        mainStack.push(x);
        if (minStack.empty()) minStack.push(x);
        else minStack.push(min(x, minStack.top()));
    }

    void pop() {
        mainStack.pop();
        minStack.pop();
    }

    int top() {
        return mainStack.top();
    }

    int getMin() {
        return minStack.top();
    }
};

int main() {
    MinStack s;
    s.push(5); s.push(3); s.push(7); s.push(2);
    s.pop();
    cout << "Min: " << s.getMin() << endl;
    return 0;
}
\end{lstlisting}

\begin{complexitybox}
\textbf{Time Complexity.} $O(1)$ for every operation.
\textbf{Space Complexity.} $O(n)$, doubled from a plain stack since every
push adds one entry to \emph{both} stacks. This can be reduced to
$O(1)$ \emph{extra} space (still $O(n)$ total, but without a second full
stack) using a trick where only the previous minimum is pushed onto the
main stack itself when a new minimum is encountered, encoded so it can be
distinguished and restored on pop -- a worthwhile optimization to know
exists, though the two-stack version above is the version to derive and
explain first.
\end{complexitybox}

\begin{takeawaybox}
Whenever an aggregate (min, max, or anything else that a plain running
variable cannot correctly \emph{revert} after a removal) must stay valid
across both insertions and removals, keep a parallel history stack rather
than trying to recompute the aggregate from scratch on every pop.
\end{takeawaybox}

\section{Valid Parentheses}
\label{sec:sq-valid-parentheses}

\problemstatement{Given a string $s$ consisting only of the characters
\texttt{'('}, \texttt{')'}, \texttt{'['}, \texttt{']'}, \texttt{'\{'},
\texttt{'\}'}, determine whether the brackets are \textbf{valid}: every
closing bracket must match the most recently opened, still-unmatched
bracket of the same type, and every opening bracket must eventually be
closed.}

\begin{patternbox}
\emph{``Matching/nesting''} is the most direct stack keyword there is: the
defining feature of valid nesting is that the bracket needing to be closed
\textbf{next} is always the \textbf{most recently opened} one still
unmatched -- exactly LIFO order. Whenever a problem describes validity in
terms of nested or paired structures (brackets, tags, nested function
calls), a stack of ``currently open, unmatched'' items is almost always
the right tool.
\end{patternbox}

\subsection*{Understanding the problem carefully}

Scan the string left to right. Every opening bracket is a commitment that
must be fulfilled later, in reverse order of when the commitments were
made -- the bracket opened most recently must be the first one closed,
because brackets cannot cross (that would not be ``nesting''). Pushing
each opening bracket onto a stack, and popping when we see a closing
bracket, mirrors this exactly: the stack's top is always the innermost
currently-open bracket, which is precisely the one any new closing
bracket must match.

\begin{ideabox}[Stack of Open Brackets]
For each character: if it is an opening bracket, push it. If it is a
closing bracket: the stack must be non-empty, and its top must be the
\emph{matching} opening bracket; if either check fails, the string is
invalid. Otherwise, pop the stack (this pair is now resolved). After
scanning the whole string, the string is valid if and only if the stack
is empty (every opened bracket was eventually closed).
\end{ideabox}

\subsection*{Why both checks (stack non-empty, and top matches) are necessary}

If the stack is empty when a closing bracket appears, there is no
open bracket left to match it against -- the closing bracket is
unmatched, and the string is invalid (e.g.\ \texttt{")"} alone). If the
stack's top does not match the closing bracket's type (e.g.\ the top is
\texttt{'('} but we just saw \texttt{']'}), the brackets have crossed
incorrectly (e.g.\ \texttt{"(]"}), which also invalidates the string. And
even if every individual closing bracket finds a match, any \emph{leftover}
opening brackets at the end (a non-empty stack) mean some opens were never
closed at all (e.g.\ \texttt{"(("}, missing both closing parentheses) --
this is why the final emptiness check is also required, not just the
per-character checks.

\subsection*{Algorithm}

\begin{enumerate}
  \item Initialize an empty stack.
  \item For each character $c$ in $s$:
  \begin{enumerate}
    \item If $c$ is an opening bracket: push $c$.
    \item Else ($c$ is a closing bracket): if the stack is empty, or its
          top is not the matching opening bracket for $c$, return
          \texttt{false}. Otherwise, pop the stack.
  \end{enumerate}
  \item Return whether the stack is now empty.
\end{enumerate}

\begin{algorithm}[H]
\caption{Valid Parentheses}
\begin{algorithmic}[1]
\Procedure{IsValid}{$s$}
  \State $\textit{stack} \gets$ empty
  \For{each character $c$ in $s$}
    \If{$c \in \{(, [, \{\}$}
      \State push $c$ onto \textit{stack}
    \Else
      \If{\textit{stack} is empty \textbf{or} top of \textit{stack} does not match $c$}
        \State \Return \textbf{false}
      \EndIf
      \State pop \textit{stack}
    \EndIf
  \EndFor
  \State \Return (\textit{stack} is empty)
\EndProcedure
\end{algorithmic}
\end{algorithm}

\subsection*{Dry run}

\dryrun{Take $s = \texttt{"\{[()]\}"}$.}

\begin{itemize}
  \item \texttt{'\{'}: push. Stack $=[\{]$.
  \item \texttt{'['}: push. Stack $=[\{,[]$.
  \item \texttt{'('}: push. Stack $=[\{,[,(]$.
  \item \texttt{')'}: top is \texttt{'('}, matches. Pop. Stack
        $=[\{,[]$.
  \item \texttt{']'}: top is \texttt{'['}, matches. Pop. Stack $=[\{]$.
  \item \texttt{'\}'}: top is \texttt{'\{'}, matches. Pop. Stack $=[]$.
\end{itemize}

Stack is empty at the end: \texttt{true}.

Now take $s = \texttt{"(]"}$: push \texttt{'('}; then see \texttt{']'},
top is \texttt{'('}, which does not match \texttt{']'} -- return
\texttt{false} immediately.

\subsection*{C++ implementation}

\begin{lstlisting}
#include <bits/stdc++.h>
using namespace std;

bool isValid(string s) {
    stack<char> st;
    for (char c : s) {
        if (c == '(' || c == '[' || c == '{') {
            st.push(c);
        } else {
            if (st.empty()) return false;
            char top = st.top();
            if ((c == ')' && top != '(') ||
                (c == ']' && top != '[') ||
                (c == '}' && top != '{')) {
                return false;
            }
            st.pop();
        }
    }
    return st.empty();
}

int main() {
    cout << (isValid("{[()]}") ? "true" : "false") << endl;
    cout << (isValid("(]") ? "true" : "false") << endl;
    return 0;
}
\end{lstlisting}

\begin{complexitybox}
\textbf{Time Complexity.} Each character is processed once with constant
work, giving
\[
  O(n).
\]
\textbf{Space Complexity.} $O(n)$ in the worst case (a string of all
opening brackets pushes every character onto the stack).
\end{complexitybox}

\begin{takeawaybox}
``The next thing that needs to be resolved is always the most recently
opened, still-unresolved thing'' is the defining signature of a stack
problem. This is not limited to brackets -- the same template applies to
validating nested tags, matching function call/return pairs, or undoing
actions in reverse order of when they were performed.
\end{takeawaybox}

\section{Infix to Postfix Conversion}
\label{sec:sq-infix-to-postfix}

\noindent We now turn to a group of six closely related problems:
converting between the three ways of writing an arithmetic expression.
\textbf{Infix} places operators between operands ($a+b$) and needs
parentheses and precedence rules to remove ambiguity. \textbf{Postfix}
(operators after their operands, $ab+$) and \textbf{prefix} (operators
before, $+ab$) are both unambiguous without any parentheses, which is
exactly why calculators and compilers prefer them internally -- and why a
stack-based scan can evaluate or convert them in a single pass.

\problemstatement{Given an infix expression (operands, the binary
operators \texttt{+ - * / \textasciicircum}, and parentheses), convert it
to postfix notation.}

\begin{patternbox}
\emph{``Convert an expression respecting operator precedence and
parentheses''} is the Shunting-Yard algorithm's signature: a stack holds
operators \emph{waiting} to be placed into the output, and the central
decision at every operator is \emph{``do any waiting operators have
\textbf{equal or higher} precedence than the one I just read, and if so,
should they come out first?''} -- exactly the kind of ordering decision a
stack resolves naturally.
\end{patternbox}

\subsection*{Understanding the problem carefully}

In postfix, an operator appears \emph{immediately after} both of its
operands are fully written -- which means an operator can only be emitted
to the output once we are certain nothing of higher (or equal, for
left-associative operators) precedence is still waiting to bind tighter
around it. Parentheses simply force this resolution explicitly: everything
inside a $(\ldots)$ must be fully resolved into postfix before the
parenthesis is allowed to close.

\begin{ideabox}[Precedence and Associativity]
Assign precedence: \texttt{\textasciicircum} (highest, right-associative)
$>$ \texttt{* /} (left-associative) $>$ \texttt{+ -} (lowest,
left-associative). When deciding whether to pop an operator already on
the stack before pushing a new one: pop if the stacked operator has
\textbf{strictly higher} precedence, or \textbf{equal} precedence \emph{and}
the new operator is left-associative (for a right-associative operator
like \texttt{\textasciicircum}, equal precedence does \emph{not} trigger a
pop, since \texttt{\textasciicircum} groups right-to-left:
$2\texttt{\textasciicircum}3\texttt{\textasciicircum}2$ means
$2\texttt{\textasciicircum}(3\texttt{\textasciicircum}2)$, not
$(2\texttt{\textasciicircum}3)\texttt{\textasciicircum}2$).
\end{ideabox}

\subsection*{Why popping by precedence produces correct postfix}

An operator should only be written to postfix output once everything that
needs to bind \emph{more tightly} around its operands has already been
resolved. If the stack's top has strictly higher precedence, it must
finish binding its own operands first (it is ``more urgent''), so we pop
it to the output before pushing the new, lower-precedence operator. If
precedences are equal and the operator is left-associative, the earlier
(already-stacked) operator should also be resolved first, since
left-associativity means we group left to right -- e.g.\ $a-b-c$ means
$(a-b)-c$, so the first \texttt{-} must close over $a,b$ before the second
\texttt{-} is considered. Right-associative operators reverse this rule
specifically to preserve right-to-left grouping.

\subsection*{Algorithm}

\begin{enumerate}
  \item Scan left to right. For each token:
  \begin{enumerate}
    \item Operand: append to output.
    \item \texttt{'('}: push.
    \item \texttt{')'}: pop and append to output until \texttt{'('} is
          popped (discard both parentheses).
    \item Operator $op$: while the stack's top is an operator with
          higher precedence, or equal precedence with $op$
          left-associative: pop and append to output. Then push $op$.
  \end{enumerate}
  \item After the scan, pop all remaining operators to the output.
\end{enumerate}

\begin{algorithm}[H]
\caption{Infix to Postfix}
\begin{algorithmic}[1]
\Procedure{InfixToPostfix}{$s$}
  \State $\textit{output} \gets$ empty, $\textit{stack} \gets$ empty
  \For{each token $t$ in $s$}
    \If{$t$ is an operand} append $t$ to \textit{output}
    \ElsIf{$t = ($} push $t$
    \ElsIf{$t = )$}
      \While{top $\ne$ (} pop, append to \textit{output} \EndWhile
      \State pop the \texttt{(}
    \Else
      \While{\textit{stack} not empty \textbf{and} top has higher precedence, \textbf{or} (equal precedence \textbf{and} $t$ is left-assoc)}
        \State pop, append to \textit{output}
      \EndWhile
      \State push $t$
    \EndIf
  \EndFor
  \While{\textit{stack} not empty} pop, append to \textit{output} \EndWhile
  \State \Return \textit{output}
\EndProcedure
\end{algorithmic}
\end{algorithm}

\subsection*{Dry run}

\dryrun{Take $s = \texttt{"a+b*c"}$.}

\begin{itemize}
  \item \texttt{a}: output $=\texttt{"a"}$.
  \item \texttt{+}: stack empty, push. Stack $=[+]$.
  \item \texttt{b}: output $=\texttt{"ab"}$.
  \item \texttt{*}: top is \texttt{+} (lower precedence than \texttt{*}),
        no pop. Push. Stack $=[+,*]$.
  \item \texttt{c}: output $=\texttt{"abc"}$.
  \item End of scan: pop all. Pop \texttt{*}: output
        $=\texttt{"abc*"}$. Pop \texttt{+}: output $=\texttt{"abc*+"}$.
\end{itemize}

Result: \texttt{"abc*+"}, correctly reflecting that \texttt{*} binds
tighter than \texttt{+} (i.e.\ $a+(b*c)$).

\subsection*{C++ implementation}

\begin{lstlisting}
#include <bits/stdc++.h>
using namespace std;

int precedence(char op) {
    if (op == '^') return 3;
    if (op == '*' || op == '/') return 2;
    if (op == '+' || op == '-') return 1;
    return 0;
}

bool isRightAssoc(char op) { return op == '^'; }

string infixToPostfix(string s) {
    string output;
    stack<char> st;

    for (char c : s) {
        if (isalnum(c)) {
            output += c;
        } else if (c == '(') {
            st.push(c);
        } else if (c == ')') {
            while (st.top() != '(') {
                output += st.top();
                st.pop();
            }
            st.pop(); // remove '('
        } else { // operator
            while (!st.empty() && st.top() != '(' &&
                   (precedence(st.top()) > precedence(c) ||
                    (precedence(st.top()) == precedence(c) && !isRightAssoc(c)))) {
                output += st.top();
                st.pop();
            }
            st.push(c);
        }
    }
    while (!st.empty()) {
        output += st.top();
        st.pop();
    }
    return output;
}

int main() {
    cout << infixToPostfix("a+b*c") << endl;
    return 0;
}
\end{lstlisting}

\begin{complexitybox}
\textbf{Time Complexity.} Each token is pushed and popped at most once,
giving
\[
  O(n).
\]
\textbf{Space Complexity.} $O(n)$ for the stack and output in the worst
case.
\end{complexitybox}

\begin{takeawaybox}
The precedence-and-associativity comparison rule derived here is reused,
completely unchanged, in Infix to Prefix
(Section~\ref{sec:sq-infix-to-prefix}) -- only the direction of scanning
and the meaning of the output change.
\end{takeawaybox}

\section{Infix to Prefix Conversion}
\label{sec:sq-infix-to-prefix}

\problemstatement{Given an infix expression, convert it to prefix
notation.}

\begin{patternbox}
\emph{``Convert to prefix''} immediately after having just solved
``convert to postfix'' (Section~\ref{sec:sq-infix-to-postfix}) is an
invitation to \textbf{reuse} that algorithm via a symmetry trick, rather
than deriving a new one: prefix is postfix \emph{read backward}, with
operators and operands swapped in scanning direction -- so reversing the
input, adapting the tie-breaking rule, running the same postfix algorithm,
and reversing the output again gets us there with no new derivation.
\end{patternbox}

\subsection*{Understanding the problem carefully}

If postfix writes operators \emph{after} their operands, prefix writes
them \emph{before} -- which is exactly what you get by mirror-imaging a
postfix expression (reversing it character by character) \emph{and}
reversing each multi-character operand if there were any. So the natural
strategy is: reverse the input expression (also swapping every
\texttt{'('} with \texttt{')'} and vice versa, since reversal flips which
parenthesis opens and which closes), run essentially the same
Shunting-Yard scan as Section~\ref{sec:sq-infix-to-postfix} on this
reversed-and-swapped string, then reverse the resulting output once more
to undo the initial reversal.

\begin{ideabox}[The One Rule That Must Change]
When running the postfix algorithm on the reversed string, the tie-break
rule for \textbf{equal precedence} must flip: do \emph{not} pop an
equal-precedence operator already on the stack before pushing a new one,
\emph{regardless of true associativity}. This is because reversing the
scan direction also reverses which operand of a left-associative chain we
encounter first; treating every tie as if it were right-associative during
this reversed scan is what correctly preserves the original
left-to-right grouping once everything is reversed back.
\end{ideabox}

\subsection*{Why reverse-scan-reverse is correct}

This is a direct application of the same ``reverse twice to flip an order
without breaking internal structure'' idea seen in
Section~\ref{sec:sq-queue-using-stack}. Reversing the infix string
converts a left-to-right scan that would naturally produce postfix into
one that, after the tie-break adjustment, produces the mirror image of
postfix -- which, when reversed back to normal reading order, is exactly
prefix. The parenthesis-swapping step is necessary because reversing a
string turns every originally-opening parenthesis into a closing one
(textually) and vice versa, and the algorithm needs them to still behave
correctly as openers and closers during the reversed scan.

\subsection*{Algorithm}

\begin{enumerate}
  \item Reverse $s$; swap every \texttt{'('} with \texttt{')'} and vice
        versa.
  \item Run the Shunting-Yard postfix algorithm
        (Algorithm~\ref{sec:sq-infix-to-postfix}) on this modified
        string, with the tie-break rule changed to \emph{never} pop on
        equal precedence.
  \item Reverse the resulting output string. This is the prefix
        expression.
\end{enumerate}

\subsection*{Dry run}

\dryrun{Take $s = \texttt{"a+b*c"}$.}

\begin{itemize}
  \item Reverse and swap brackets (no brackets here): $\texttt{"c*b+a"}$.
  \item Run the postfix algorithm (never popping on ties):
  \begin{itemize}
    \item \texttt{c}: output $=\texttt{"c"}$.
    \item \texttt{*}: push. Stack $=[*]$.
    \item \texttt{b}: output $=\texttt{"cb"}$.
    \item \texttt{+}: top is \texttt{*} (higher precedence), pop:
          output $=\texttt{"cb*"}$. Push \texttt{+}. Stack $=[+]$.
    \item \texttt{a}: output $=\texttt{"cb*a"}$.
    \item End: pop \texttt{+}: output $=\texttt{"cb*a+"}$.
  \end{itemize}
  \item Reverse the output: $\texttt{"+a*bc"}$.
\end{itemize}

Result: \texttt{"+a*bc"}, correctly representing $a+(b*c)$ in prefix.

\subsection*{C++ implementation}

\begin{lstlisting}
#include <bits/stdc++.h>
using namespace std;

int precedence(char op) {
    if (op == '^') return 3;
    if (op == '*' || op == '/') return 2;
    if (op == '+' || op == '-') return 1;
    return 0;
}

string infixToPrefix(string s) {
    reverse(s.begin(), s.end());
    for (char& c : s) {
        if (c == '(') c = ')';
        else if (c == ')') c = '(';
    }

    string output;
    stack<char> st;

    for (char c : s) {
        if (isalnum(c)) {
            output += c;
        } else if (c == '(') {
            st.push(c);
        } else if (c == ')') {
            while (st.top() != '(') {
                output += st.top();
                st.pop();
            }
            st.pop();
        } else { // operator: never pop on equal precedence
            while (!st.empty() && st.top() != '(' &&
                   precedence(st.top()) > precedence(c)) {
                output += st.top();
                st.pop();
            }
            st.push(c);
        }
    }
    while (!st.empty()) {
        output += st.top();
        st.pop();
    }

    reverse(output.begin(), output.end());
    return output;
}

int main() {
    cout << infixToPrefix("a+b*c") << endl;
    return 0;
}
\end{lstlisting}

\begin{complexitybox}
\textbf{Time Complexity.} $O(n)$, identical in structure to
Section~\ref{sec:sq-infix-to-postfix}.
\textbf{Space Complexity.} $O(n)$.
\end{complexitybox}

\begin{takeawaybox}
``Reverse, adapt the tie-break rule, run the algorithm you already have,
reverse again'' is a powerful general technique for converting a
left-to-right algorithm into its mirror-image counterpart, without
re-deriving the core logic. Recognize when a new problem is the
\emph{mirror image} of one you have already solved before writing new
code from scratch.
\end{takeawaybox}

\section{Postfix to Infix Conversion}
\label{sec:sq-postfix-to-infix}

\problemstatement{Given a postfix expression, convert it to infix
notation (fully parenthesized, to remove any ambiguity).}

\begin{patternbox}
Converting \emph{from} postfix or prefix is structurally easier than
converting \emph{from} infix (Sections~\ref{sec:sq-infix-to-postfix}--\ref{sec:sq-infix-to-prefix}),
because postfix and prefix already encode the expression tree's structure
unambiguously through ordering alone -- there is no precedence or
associativity to resolve. The keyword shift to notice: when the
\emph{source} notation is postfix or prefix, the stack holds
\textbf{partially built sub-expressions} (strings), not raw operators, and
every operator immediately and unambiguously tells you which two
sub-expressions to combine.
\end{patternbox}

\subsection*{Understanding the problem carefully}

In postfix, by the time we encounter an operator while scanning left to
right, both of its operands have already been fully read and are sitting
on top of the stack (the most recent two pushes) -- there is never any
ambiguity about which sub-expressions an operator applies to, unlike
infix where precedence must be consulted.

\begin{ideabox}[Stack of Sub-expressions]
Scan left to right. Push every operand as a (string) sub-expression. On
an operator: pop the top two sub-expressions -- call the first popped
$\textit{right}$ and the second popped $\textit{left}$ (since
$\textit{left}$ was pushed earlier and so sits one level deeper). Combine
them as $(\textit{left}\ \textit{op}\ \textit{right})$, fully
parenthesized, and push this combined string back as a single
sub-expression.
\end{ideabox}

\subsection*{Why the most recent two pushes are always the correct operands}

Postfix notation guarantees that an operator always immediately follows
the complete sub-expressions it operates on -- this is precisely what
``postfix'' (operator \emph{after} operands) means. So whatever is on top
of the stack at the moment we read an operator must be its right operand
(read most recently), and whatever is just below that must be its left
operand, with no possibility of some other, older sub-expression
interfering, because every operator so far has already fully consumed and
replaced its own operands before we reach this point.

\subsection*{Algorithm}

\begin{enumerate}
  \item Initialize an empty stack of strings.
  \item For each token in the postfix expression:
  \begin{enumerate}
    \item Operand: push it as a one-character string.
    \item Operator $\textit{op}$: pop $\textit{right}$, pop
          $\textit{left}$; push $(\textit{left}\ \textit{op}\
          \textit{right})$.
  \end{enumerate}
  \item The stack's single remaining element is the infix result.
\end{enumerate}

\subsection*{Dry run}

\dryrun{Take $\texttt{"ab+c*"}$.}

\begin{itemize}
  \item \texttt{a}: push. Stack $=[\texttt{a}]$.
  \item \texttt{b}: push. Stack $=[\texttt{a},\texttt{b}]$.
  \item \texttt{+}: pop right$=\texttt{b}$, pop left$=\texttt{a}$. Push
        \texttt{(a+b)}. Stack $=[\texttt{(a+b)}]$.
  \item \texttt{c}: push. Stack $=[\texttt{(a+b)},\texttt{c}]$.
  \item \texttt{*}: pop right$=\texttt{c}$, pop left$=\texttt{(a+b)}$.
        Push \texttt{((a+b)*c)}.
\end{itemize}

Result: \texttt{((a+b)*c)}.

\subsection*{C++ implementation}

\begin{lstlisting}
#include <bits/stdc++.h>
using namespace std;

string postfixToInfix(string s) {
    stack<string> st;

    for (char c : s) {
        if (isalnum(c)) {
            st.push(string(1, c));
        } else {
            string right = st.top(); st.pop();
            string left = st.top(); st.pop();
            st.push("(" + left + string(1, c) + right + ")");
        }
    }
    return st.top();
}

int main() {
    cout << postfixToInfix("ab+c*") << endl;
    return 0;
}
\end{lstlisting}

\begin{complexitybox}
\textbf{Time Complexity.} Each token triggers $O(1)$ stack operations,
but string concatenation can cost up to $O(n)$ per operator in the worst
case (building progressively longer strings), giving an overall worst
case of
\[
  O(n^2)
\]
for naive string concatenation (this can be improved to $O(n)$ with more
careful string-building techniques, but $O(n^2)$ is what the straightforward
implementation above gives).
\textbf{Space Complexity.} $O(n)$ for the stack of sub-expressions.
\end{complexitybox}

\begin{takeawaybox}
When converting \emph{from} postfix or prefix, the stack holds
sub-expressions (strings), not bare operators -- every operator
encountered has an unambiguous, immediate pair of operands to combine,
with no precedence comparison ever needed. This is structurally simpler
than converting \emph{from} infix, and the same ``stack of
sub-expressions'' idea reused, with only the combination format
(parenthesized infix here, vs.\ prefix or postfix in the next few
sections) changing.
\end{takeawaybox}

\section{Postfix to Prefix Conversion}
\label{sec:sq-postfix-to-prefix}

\problemstatement{Given a postfix expression, convert it to prefix
notation.}

\begin{patternbox}
Identical setup to Section~\ref{sec:sq-postfix-to-infix} -- the only
change is \emph{how} two popped sub-expressions are combined: instead of
$(\textit{left}\ \textit{op}\ \textit{right})$, prefix writes the operator
\textbf{first}, with no parentheses needed at all (prefix, like postfix,
is unambiguous without them).
\end{patternbox}

\subsection*{Understanding the problem carefully}

The stack-of-sub-expressions technique from
Section~\ref{sec:sq-postfix-to-infix} carries over completely unchanged;
only the string-building step at each operator differs.

\begin{ideabox}[Combine as Prefix Instead of Infix]
On an operator $\textit{op}$: pop $\textit{right}$, pop $\textit{left}$
(same order as before); push $\textit{op} + \textit{left} +
\textit{right}$ (operator first, then both sub-expressions
back-to-back, no parentheses or spaces needed since prefix length is
self-delimiting when operands are single characters, or separated by a
delimiter for multi-character operands).
\end{ideabox}

\subsection*{Algorithm}

\begin{enumerate}
  \item Initialize an empty stack of strings.
  \item For each token: operand $\to$ push as a string. Operator
        $\textit{op}$ $\to$ pop $\textit{right}$, pop $\textit{left}$;
        push $\textit{op} + \textit{left} + \textit{right}$.
  \item The stack's single remaining element is the prefix result.
\end{enumerate}

\subsection*{Dry run}

\dryrun{Take $\texttt{"ab+c*"}$.}

\begin{itemize}
  \item \texttt{a}: push. Stack $=[\texttt{a}]$.
  \item \texttt{b}: push. Stack $=[\texttt{a},\texttt{b}]$.
  \item \texttt{+}: pop right$=\texttt{b}$, left$=\texttt{a}$. Push
        \texttt{+ab}. Stack $=[\texttt{+ab}]$.
  \item \texttt{c}: push. Stack $=[\texttt{+ab},\texttt{c}]$.
  \item \texttt{*}: pop right$=\texttt{c}$, left$=\texttt{+ab}$. Push
        \texttt{*+abc}.
\end{itemize}

Result: \texttt{*+abc}, correctly representing $(a+b)*c$ in prefix.

\subsection*{C++ implementation}

\begin{lstlisting}
#include <bits/stdc++.h>
using namespace std;

string postfixToPrefix(string s) {
    stack<string> st;

    for (char c : s) {
        if (isalnum(c)) {
            st.push(string(1, c));
        } else {
            string right = st.top(); st.pop();
            string left = st.top(); st.pop();
            st.push(string(1, c) + left + right);
        }
    }
    return st.top();
}

int main() {
    cout << postfixToPrefix("ab+c*") << endl;
    return 0;
}
\end{lstlisting}

\begin{complexitybox}
\textbf{Time Complexity.} $O(n^2)$ worst case with naive string
concatenation, as in Section~\ref{sec:sq-postfix-to-infix}.
\textbf{Space Complexity.} $O(n)$.
\end{complexitybox}

\begin{takeawaybox}
This section is intentionally a near-copy of
Section~\ref{sec:sq-postfix-to-infix} -- the lesson is that once you have
the ``stack of sub-expressions, combine on every operator'' skeleton
right, every output \emph{format} (parenthesized infix, prefix, or
postfix-of-a-different-source) is just a one-line change to how the
combination step builds its string.
\end{takeawaybox}

\section{Prefix to Infix Conversion}
\label{sec:sq-prefix-to-infix}

\problemstatement{Given a prefix expression, convert it to infix notation
(fully parenthesized).}

\begin{patternbox}
Prefix writes the operator \emph{before} its operands, so by the time we
encounter an operator while scanning \textbf{left to right}, its operands
have not been read yet -- the stack-of-sub-expressions trick from
Sections~\ref{sec:sq-postfix-to-infix}--\ref{sec:sq-postfix-to-prefix}
only works smoothly when operands are already available \emph{before}
their operator is seen. The fix: scan \textbf{right to left} instead,
which makes prefix behave exactly like postfix did under a left-to-right
scan.
\end{patternbox}

\subsection*{Understanding the problem carefully}

Reading a prefix expression from right to left, an operand is encountered
\emph{before} the operator that uses it (mirroring how, in postfix read
left to right, an operand is also encountered before its operator) -- so
the same stack-of-sub-expressions technique applies, just with the scan
direction reversed.

\begin{ideabox}[Scan Right to Left, Same Combination Rule]
Scan from the last character to the first. Push every operand. On an
operator $\textit{op}$: pop the top sub-expression as $\textit{left}$
(it was pushed most recently, and since we are scanning backward, ``most
recent'' corresponds to the operand that appears \emph{first}, i.e.\
leftmost, after this operator when read normally); pop the next as
$\textit{right}$. Push $(\textit{left}\ \textit{op}\ \textit{right})$.
\end{ideabox}

\subsection*{Why the pop order flips relative to postfix}

In postfix, scanning left to right, the operand pushed \emph{last}
(closest to the operator, on its left in the text) is the \textbf{right}
operand semantically (e.g.\ in $\texttt{ab+}$, $b$ is pushed right before
\texttt{+}, and $a+b$ has $b$ as the right operand). In prefix, scanning
\emph{right to left}, the operand encountered (and pushed) most recently
before the operator is the one immediately to the operator's right in the
original text -- which, in prefix's $\textit{op}\ A\ B$ structure, is $A$,
the \textbf{left} operand. So the pop order is mirrored relative to
postfix's, which is exactly why $\textit{left}$ is popped first here.

\subsection*{Algorithm}

\begin{enumerate}
  \item Initialize an empty stack of strings.
  \item For each token, scanning from right to left:
  \begin{enumerate}
    \item Operand: push it.
    \item Operator $\textit{op}$: pop $\textit{left}$, pop
          $\textit{right}$; push $(\textit{left}\ \textit{op}\
          \textit{right})$.
  \end{enumerate}
  \item The stack's single remaining element is the infix result.
\end{enumerate}

\subsection*{Dry run}

\dryrun{Take $\texttt{"*+abc"}$ (representing $(a+b)*c$).}

Scanning right to left: \texttt{c}, \texttt{b}, \texttt{a}, \texttt{+},
\texttt{*}.

\begin{itemize}
  \item \texttt{c}: push. Stack $=[\texttt{c}]$.
  \item \texttt{b}: push. Stack $=[\texttt{c},\texttt{b}]$.
  \item \texttt{a}: push. Stack $=[\texttt{c},\texttt{b},\texttt{a}]$.
  \item \texttt{+}: pop left$=\texttt{a}$, pop right$=\texttt{b}$. Push
        \texttt{(a+b)}. Stack $=[\texttt{c},\texttt{(a+b)}]$.
  \item \texttt{*}: pop left$=\texttt{(a+b)}$, pop right$=\texttt{c}$.
        Push \texttt{((a+b)*c)}.
\end{itemize}

Result: \texttt{((a+b)*c)} -- matching
Section~\ref{sec:sq-postfix-to-infix}'s result for the same underlying
expression.

\subsection*{C++ implementation}

\begin{lstlisting}
#include <bits/stdc++.h>
using namespace std;

string prefixToInfix(string s) {
    stack<string> st;

    for (int i = s.size() - 1; i >= 0; i--) {
        char c = s[i];
        if (isalnum(c)) {
            st.push(string(1, c));
        } else {
            string left = st.top(); st.pop();
            string right = st.top(); st.pop();
            st.push("(" + left + string(1, c) + right + ")");
        }
    }
    return st.top();
}

int main() {
    cout << prefixToInfix("*+abc") << endl;
    return 0;
}
\end{lstlisting}

\begin{complexitybox}
\textbf{Time Complexity.} $O(n^2)$ worst case with naive string
concatenation, $O(n)$ with more careful string building.
\textbf{Space Complexity.} $O(n)$.
\end{complexitybox}

\begin{takeawaybox}
Scanning right to left is to prefix what scanning left to right is to
postfix -- both make operands available before the operator that needs
them. Whenever a problem's natural scan direction would otherwise present
an operator before its operands are ready, check whether reversing the
scan direction restores the same simple stack technique.
\end{takeawaybox}

\section{Prefix to Postfix Conversion}
\label{sec:sq-prefix-to-postfix}

\problemstatement{Given a prefix expression, convert it to postfix
notation.}

\begin{patternbox}
The last of the six conversions, and by now a pure exercise in combining
two already-derived ideas: scan \textbf{right to left} (as in
Section~\ref{sec:sq-prefix-to-infix}, since the source is prefix), and
combine each pair as $\textit{left} + \textit{right} + \textit{op}$ (as in
Section~\ref{sec:sq-postfix-to-prefix}'s combination style, just operator
\emph{last} instead of first, and no parentheses needed).
\end{patternbox}

\subsection*{Understanding the problem carefully}

No new idea is required: take the right-to-left scan direction from
Section~\ref{sec:sq-prefix-to-infix} (because the source is prefix) and
the parenthesis-free, operator-position-only combination style appropriate
for postfix output (operator written last).

\begin{ideabox}[Combine as Postfix Instead of Infix]
Scan right to left. Push every operand. On operator $\textit{op}$: pop
$\textit{left}$, pop $\textit{right}$ (same pop order as
Section~\ref{sec:sq-prefix-to-infix}, since the source notation and scan
direction are unchanged); push $\textit{left} + \textit{right} +
\textit{op}$.
\end{ideabox}

\subsection*{Algorithm}

\begin{enumerate}
  \item Initialize an empty stack of strings.
  \item For each token, scanning right to left: operand $\to$ push.
        Operator $\textit{op}$ $\to$ pop $\textit{left}$, pop
        $\textit{right}$; push $\textit{left} + \textit{right} +
        \textit{op}$.
  \item The stack's single remaining element is the postfix result.
\end{enumerate}

\subsection*{Dry run}

\dryrun{Take $\texttt{"*+abc"}$.}

Scanning right to left: \texttt{c}, \texttt{b}, \texttt{a}, \texttt{+},
\texttt{*}.

\begin{itemize}
  \item \texttt{c}: push. Stack $=[\texttt{c}]$.
  \item \texttt{b}: push. Stack $=[\texttt{c},\texttt{b}]$.
  \item \texttt{a}: push. Stack $=[\texttt{c},\texttt{b},\texttt{a}]$.
  \item \texttt{+}: pop left$=\texttt{a}$, pop right$=\texttt{b}$. Push
        \texttt{ab+}. Stack $=[\texttt{c},\texttt{ab+}]$.
  \item \texttt{*}: pop left$=\texttt{ab+}$, pop right$=\texttt{c}$.
        Push \texttt{ab+c*}.
\end{itemize}

Result: \texttt{ab+c*} -- exactly the postfix expression we started from
back in Section~\ref{sec:sq-postfix-to-infix}, confirming the whole
six-conversion group is internally consistent.

\subsection*{C++ implementation}

\begin{lstlisting}
#include <bits/stdc++.h>
using namespace std;

string prefixToPostfix(string s) {
    stack<string> st;

    for (int i = s.size() - 1; i >= 0; i--) {
        char c = s[i];
        if (isalnum(c)) {
            st.push(string(1, c));
        } else {
            string left = st.top(); st.pop();
            string right = st.top(); st.pop();
            st.push(left + right + string(1, c));
        }
    }
    return st.top();
}

int main() {
    cout << prefixToPostfix("*+abc") << endl;
    return 0;
}
\end{lstlisting}

\begin{complexitybox}
\textbf{Time Complexity.} $O(n^2)$ worst case with naive string
concatenation, $O(n)$ with more careful string building.
\textbf{Space Complexity.} $O(n)$.
\end{complexitybox}

\begin{takeawaybox}
Across all six conversions in this group, there are really only two
independent decisions: \emph{which direction to scan} (left to right if
the source is postfix or plain infix; right to left if the source is
prefix) and \emph{how to format the combination step} (parenthesized
infix, operator-first prefix, or operator-last postfix). Every one of the
six problems is a different combination of these same two choices.
\end{takeawaybox}

\section{Next Greater Element I}
\label{sec:sq-next-greater-1}

\noindent We now turn to the second major group of problems in this
chapter: the \textbf{monotonic stack}, introduced in
Section~\ref{sec:sq-intro}. This first problem is the purest possible
instance of the pattern, and every later monotonic-stack problem in this
chapter is a variation or extension of the technique derived here.

\problemstatement{We are given two arrays, $\textit{nums1}$ and
$\textit{nums2}$, both containing distinct integers, where every element
of $\textit{nums1}$ is also present somewhere in $\textit{nums2}$. For
each element of $\textit{nums1}$, find its \textbf{next greater element}
in $\textit{nums2}$: the first element to its right, within
$\textit{nums2}$, that is strictly greater than it. If no such element
exists, the answer for that element is $-1$.}

\begin{patternbox}
\emph{``Next greater element''} is, almost word for word, the keyword
from Section~\ref{sec:sq-intro}'s pattern list. The structure here is
two-part: first solve ``next greater element, for every position, within
\textit{nums2}'' (a self-contained monotonic-stack problem), then a simple
hash-map lookup translates that solved array into answers for
$\textit{nums1}$.
\end{patternbox}

\subsection*{Understanding the problem carefully}

The core sub-problem is: for every index $i$ in $\textit{nums2}$, find the
nearest index $j>i$ with $\textit{nums2}[j] > \textit{nums2}[i]$, or
report none exists. A brute-force solution checks every pair, costing
$O(n^2)$. The monotonic stack from Section~\ref{sec:sq-intro} solves this
in $O(n)$ by scanning right to left and maintaining a stack that is always
strictly decreasing from bottom to top.

\begin{ideabox}[Monotonic Decreasing Stack, Scanned Right to Left]
Scan $\textit{nums2}$ from right to left. Before processing
$\textit{nums2}[i]$, pop every stack element $\le \textit{nums2}[i]$
(they can never be the answer for anything still to the left, since
$\textit{nums2}[i]$ is closer and at least as large). Whatever remains on
top of the stack after this popping (if anything) is exactly
$i$'s next greater element. Then push $\textit{nums2}[i]$.
\end{ideabox}

\subsection*{Why this gives the correct next greater element for every index}

After popping every $\le \textit{nums2}[i]$ value, the stack's remaining
top -- if any -- is the \emph{nearest} value to the right of $i$ that
exceeds $\textit{nums2}[i]$. This holds because the stack only ever
contains values in decreasing order from bottom to top, and every value
currently on the stack represents some index to the right of $i$ that has
not yet been ``defeated'' by anything closer and at least as large; the
topmost (most recently pushed, hence nearest) surviving value is therefore
guaranteed to be the closest qualifying greater value. Once computed for
every index of $\textit{nums2}$, store the results in a hash map keyed by
value (values are distinct, so this is unambiguous), and answer each
query for $\textit{nums1}$ with a single $O(1)$ lookup.

\subsection*{Algorithm}

\begin{enumerate}
  \item Initialize an empty stack and an empty hash map
        $\textit{nextGreater}$.
  \item For $i$ from $n-1$ down to $0$ (scanning $\textit{nums2}$ right
        to left):
  \begin{enumerate}
    \item While the stack is not empty and its top $\le
          \textit{nums2}[i]$: pop.
    \item $\textit{nextGreater}[\textit{nums2}[i]] \gets$ the stack's top
          if non-empty, else $-1$.
    \item Push $\textit{nums2}[i]$.
  \end{enumerate}
  \item For each query value in $\textit{nums1}$: look up
        $\textit{nextGreater}[\textit{value}]$.
\end{enumerate}

\begin{algorithm}[H]
\caption{Next Greater Element I}
\begin{algorithmic}[1]
\Procedure{NextGreaterElement}{\textit{nums1}, \textit{nums2}}
  \State $\textit{stack} \gets$ empty, $\textit{nextGreater} \gets$ empty map
  \For{$i \gets n-1$ \textbf{down to} $0$}
    \While{\textit{stack} not empty \textbf{and} top $\le \textit{nums2}[i]$}
      \State pop
    \EndWhile
    \State $\textit{nextGreater}[\textit{nums2}[i]] \gets$ (\textit{stack} empty $? -1 :$ top)
    \State push $\textit{nums2}[i]$
  \EndFor
  \State \Return $[\textit{nextGreater}[x]$ for each $x$ in $\textit{nums1}]$
\EndProcedure
\end{algorithmic}
\end{algorithm}

\subsection*{Dry run}

\dryrun{Take $\textit{nums2} = [1,3,4,2]$, $\textit{nums1} = [4,1,2]$.}

Scanning right to left:

\begin{itemize}
  \item $i=3$ (value $2$): stack empty. $\textit{nextGreater}[2]=-1$.
        Push $2$. Stack $=[2]$.
  \item $i=2$ (value $4$): top is $2 \le 4$, pop. Stack empty.
        $\textit{nextGreater}[4]=-1$. Push $4$. Stack $=[4]$.
  \item $i=1$ (value $3$): top is $4 > 3$, no pop.
        $\textit{nextGreater}[3]=4$. Push $3$. Stack $=[4,3]$.
  \item $i=0$ (value $1$): top is $3 > 1$, no pop.
        $\textit{nextGreater}[1]=3$. Push $1$. Stack $=[4,3,1]$.
\end{itemize}

Map: $\{2{:}{-1}, 4{:}{-1}, 3{:}4, 1{:}3\}$. Querying $\textit{nums1} =
[4,1,2]$: answers $[-1, 3, -1]$.

\subsection*{C++ implementation}

\begin{lstlisting}
#include <bits/stdc++.h>
using namespace std;

vector<int> nextGreaterElement(vector<int>& nums1, vector<int>& nums2) {
    unordered_map<int,int> nextGreater;
    stack<int> st;

    for (int i = nums2.size() - 1; i >= 0; i--) {
        while (!st.empty() && st.top() <= nums2[i]) {
            st.pop();
        }
        nextGreater[nums2[i]] = st.empty() ? -1 : st.top();
        st.push(nums2[i]);
    }

    vector<int> result;
    for (int x : nums1) result.push_back(nextGreater[x]);
    return result;
}

int main() {
    vector<int> nums2 = {1, 3, 4, 2};
    vector<int> nums1 = {4, 1, 2};
    for (int x : nextGreaterElement(nums1, nums2)) cout << x << " ";
    cout << endl;
    return 0;
}
\end{lstlisting}

\begin{complexitybox}
\textbf{Time Complexity.} Each element of $\textit{nums2}$ is pushed once
and popped at most once across the entire scan, giving $O(n)$ for the
preprocessing; each of the $m$ queries is an $O(1)$ map lookup, giving an
overall time complexity of
\[
  O(n + m).
\]
\textbf{Space Complexity.} $O(n)$ for the stack and hash map.
\end{complexitybox}

\begin{takeawaybox}
The monotonic stack derived here -- scan right to left, pop everything
$\le$ the current element, the survivor on top is the answer, then push
-- is the master template for this entire sub-chapter. Every remaining
monotonic-stack section either applies this template directly (next
smaller, previous greater/smaller) or builds one more layer on top of it
(histogram area, trapping rain water, subarray sums).
\end{takeawaybox}

\section{Next Greater Element II (Circular Array)}
\label{sec:sq-next-greater-2}

\problemstatement{We are given a \textbf{circular} array $\textit{nums}$
(the element after the last one wraps around to the first). For every
element, find its next greater element, searching circularly -- i.e.\ we
may wrap around the end of the array once, but not more than once.}

\begin{patternbox}
The keyword \emph{``circular''} attached to a next-greater-element problem
signals that the search for each element's answer may need to continue
past the array's literal end and wrap to the beginning -- but only
\textbf{once around}, not indefinitely. The standard trick for simulating
one extra lap around a circular array without actually duplicating it in
memory is to iterate over \textbf{twice} the indices ($0$ to $2n-1$), using
the modulo operation to map each iteration index back into the real array.
\end{patternbox}

\subsection*{Understanding the problem carefully}

If we simply ran the Section~\ref{sec:sq-next-greater-1} algorithm once
over $\textit{nums}$, an element near the \emph{end} of the array would
never get a chance to find a next-greater candidate that happens to sit
near the \emph{beginning} -- the algorithm would incorrectly report $-1$
for it. Conceptually, we want to pretend the array is duplicated end to
end ($\textit{nums} + \textit{nums}$) and search within that doubled view,
but stop considering circular wraparound after one full lap (we must not
wrap twice, or an element could incorrectly become its own next greater
element).

\begin{ideabox}[Simulate a Doubled Array via Modulo Indexing]
Run the same right-to-left monotonic-stack scan as
Section~\ref{sec:sq-next-greater-1}, but let $i$ range from $2n-1$ down
to $0$, accessing $\textit{nums}[i \bmod n]$ at each step. Only
\emph{record} an answer for $i < n$ (the first lap, which corresponds to
real array positions); the second lap ($i \ge n$, processed first since we
scan from $2n-1$ downward) exists purely to \emph{seed the stack} with
candidates that wrap around from the beginning of the array, without
double-counting any element as its own answer.
\end{ideabox}

\subsection*{Why processing the second lap first, without recording answers, is correct}

Scanning from $2n-1$ down to $0$ means we process the conceptual ``second
copy'' of the array (indices $n$ to $2n-1$, mapping to real indices $0$
to $n-1$) \emph{before} the conceptual ``first copy'' (indices $0$ to
$n-1$). By the time we reach real index $i$ in the first-copy pass (i.e.\
when the loop variable itself is also $i$, having already processed
$i+1, \dots, 2n-1$), the stack already reflects every element that could
serve as $i$'s next greater element either later in the same lap
\emph{or} wrapped around from the start of the array (since those wrapped
candidates were pushed during the second-lap pass that ran first). This
correctly captures one full circular lookahead without ever risking a
second wraparound, since each real index contributes to the stack at most
twice (once per lap) and answers are only recorded during the final,
first-copy pass.

\subsection*{Algorithm}

\begin{enumerate}
  \item Initialize an empty stack and a result array of size $n$, filled
        with $-1$.
  \item For $i$ from $2n-1$ down to $0$:
  \begin{enumerate}
    \item Let $\textit{val} \gets \textit{nums}[i \bmod n]$.
    \item While the stack is not empty and its top $\le \textit{val}$:
          pop.
    \item If $i < n$: set $\textit{result}[i] \gets$ the stack's top if
          non-empty, else $-1$.
    \item Push \textit{val}.
  \end{enumerate}
  \item Return \textit{result}.
\end{enumerate}

\begin{algorithm}[H]
\caption{Next Greater Element II}
\begin{algorithmic}[1]
\Procedure{NextGreaterElementsCircular}{\textit{nums}}
  \State $\textit{result}[0..n-1] \gets -1$, $\textit{stack} \gets$ empty
  \For{$i \gets 2n-1$ \textbf{down to} $0$}
    \State $\textit{val} \gets \textit{nums}[i \bmod n]$
    \While{\textit{stack} not empty \textbf{and} top $\le \textit{val}$}
      \State pop
    \EndWhile
    \If{$i < n$ \textbf{and} \textit{stack} not empty}
      \State $\textit{result}[i] \gets$ top
    \EndIf
    \State push \textit{val}
  \EndFor
  \State \Return \textit{result}
\EndProcedure
\end{algorithmic}
\end{algorithm}

\subsection*{Dry run}

\dryrun{Take $\textit{nums} = [1,2,1]$ ($n=3$).}

\begin{itemize}
  \item $i=5$ ($\textit{val}=\textit{nums}[2]=1$): stack empty. $i\ge n$,
        no record. Push $1$. Stack $=[1]$.
  \item $i=4$ ($\textit{val}=\textit{nums}[1]=2$): top $1\le2$, pop.
        Stack empty. $i\ge n$, no record. Push $2$. Stack $=[2]$.
  \item $i=3$ ($\textit{val}=\textit{nums}[0]=1$): top $2>1$, no pop.
        $i\ge n$, no record. Push $1$. Stack $=[2,1]$.
  \item $i=2$ ($\textit{val}=\textit{nums}[2]=1$): top $1\le1$, pop.
        New top $2>1$, no further pop. $i<n$: $\textit{result}[2]=2$.
        Push $1$. Stack $=[2,1]$.
  \item $i=1$ ($\textit{val}=\textit{nums}[1]=2$): top $1\le2$, pop. New
        top $2\le2$, pop. Stack empty. $i<n$:
        $\textit{result}[1]=-1$. Push $2$. Stack $=[2]$.
  \item $i=0$ ($\textit{val}=\textit{nums}[0]=1$): top $2>1$, no pop.
        $i<n$: $\textit{result}[0]=2$. Push $1$. Stack $=[2,1]$.
\end{itemize}

Result: $[2,-1,2]$. Checking by hand: index $0$ (value $1$) wraps to find
$2$ at index $1$; index $1$ (value $2$) has no greater element anywhere
(it is the maximum); index $2$ (value $1$) wraps around to find $2$ at
index $1$ (or directly via the doubled view, finds the $2$ that appears
right after it when wrapping).

\subsection*{C++ implementation}

\begin{lstlisting}
#include <bits/stdc++.h>
using namespace std;

vector<int> nextGreaterElementsCircular(vector<int>& nums) {
    int n = nums.size();
    vector<int> result(n, -1);
    stack<int> st;

    for (int i = 2 * n - 1; i >= 0; i--) {
        int val = nums[i % n];
        while (!st.empty() && st.top() <= val) {
            st.pop();
        }
        if (i < n && !st.empty()) {
            result[i] = st.top();
        }
        st.push(val);
    }
    return result;
}

int main() {
    vector<int> nums = {1, 2, 1};
    for (int x : nextGreaterElementsCircular(nums)) cout << x << " ";
    cout << endl;
    return 0;
}
\end{lstlisting}

\begin{complexitybox}
\textbf{Time Complexity.} The loop runs $2n$ iterations, each doing
amortized $O(1)$ stack work (every element is pushed at most twice and
popped at most twice across the whole run), giving
\[
  O(n).
\]
\textbf{Space Complexity.} $O(n)$ for the stack and result array.
\end{complexitybox}

\begin{takeawaybox}
``Circular array, search wrapping around once'' has a standard, reusable
fix: iterate over $2n$ virtual indices with modulo wraparound, using the
first lap purely to seed state (here, the monotonic stack) and the second
lap to record real answers. This avoids physically duplicating the array
and generalizes to other circular-array problems beyond monotonic stacks.
\end{takeawaybox}

\section{Next Smaller Element}
\label{sec:sq-next-smaller}

\problemstatement{Given an array $\textit{arr}$, find, for every element,
the next element to its right that is strictly \textbf{smaller}, or $-1$
if none exists.}

\begin{patternbox}
The mirror image of Section~\ref{sec:sq-next-greater-1}: flip every
comparison, and the monotonic stack should now be kept strictly
\textbf{increasing} from bottom to top instead of decreasing.
\end{patternbox}

\subsection*{Understanding the problem carefully}

By the same reasoning as next greater element, scanning right to left:
before processing $\textit{arr}[i]$, any stack element $\ge
\textit{arr}[i]$ can never be the ``next smaller'' answer for anything
still to the left, since $\textit{arr}[i]$ is closer and at least as
small. Popping these and keeping whatever survives gives the nearest
qualifying smaller value.

\begin{ideabox}[Monotonic Increasing Stack, Scanned Right to Left]
Scan right to left. Before processing $\textit{arr}[i]$, pop every stack
element $\ge \textit{arr}[i]$. The survivor on top (if any) is $i$'s next
smaller element. Push $\textit{arr}[i]$.
\end{ideabox}

\subsection*{Algorithm}

\begin{enumerate}
  \item Initialize an empty stack and a result array.
  \item For $i$ from $n-1$ down to $0$:
  \begin{enumerate}
    \item While the stack is not empty and its top $\ge
          \textit{arr}[i]$: pop.
    \item $\textit{result}[i] \gets$ the stack's top if non-empty, else
          $-1$.
    \item Push $\textit{arr}[i]$.
  \end{enumerate}
  \item Return \textit{result}.
\end{enumerate}

\subsection*{Dry run}

\dryrun{Take $\textit{arr} = [4,8,5,2,25]$.}

\begin{itemize}
  \item $i=4$ ($25$): stack empty. $\textit{result}[4]=-1$. Push $25$.
  \item $i=3$ ($2$): top $25\ge2$, pop. Stack empty.
        $\textit{result}[3]=-1$. Push $2$.
  \item $i=2$ ($5$): top $2<5$, no pop. $\textit{result}[2]=2$. Push
        $5$. Stack $=[2,5]$.
  \item $i=1$ ($8$): top $5<8$, no pop. $\textit{result}[1]=5$. Push
        $8$. Stack $=[2,5,8]$.
  \item $i=0$ ($4$): top $8\ge4$, pop. New top $5\ge4$, pop. New top
        $2<4$, no pop. $\textit{result}[0]=2$. Push $4$.
\end{itemize}

Result: $[2,5,2,-1,-1]$.

\subsection*{C++ implementation}

\begin{lstlisting}
#include <bits/stdc++.h>
using namespace std;

vector<int> nextSmallerElement(vector<int>& arr) {
    int n = arr.size();
    vector<int> result(n, -1);
    stack<int> st;

    for (int i = n - 1; i >= 0; i--) {
        while (!st.empty() && st.top() >= arr[i]) {
            st.pop();
        }
        result[i] = st.empty() ? -1 : st.top();
        st.push(arr[i]);
    }
    return result;
}

int main() {
    vector<int> arr = {4, 8, 5, 2, 25};
    for (int x : nextSmallerElement(arr)) cout << x << " ";
    cout << endl;
    return 0;
}
\end{lstlisting}

\begin{complexitybox}
\textbf{Time Complexity.} $O(n)$, by the same amortized argument as
Section~\ref{sec:sq-next-greater-1}.
\textbf{Space Complexity.} $O(n)$.
\end{complexitybox}

\begin{takeawaybox}
``Next/previous greater/smaller'' is really one template with four
dials: which direction to scan (left-to-right for \emph{previous},
right-to-left for \emph{next}) and which comparison the stack enforces
(decreasing for \emph{greater}, increasing for \emph{smaller}). Once you
have derived one of the four, the rest follow by flipping the relevant
dial.
\end{takeawaybox}

\section{Number of NGEs to the Right}
\label{sec:sq-num-nges}

\problemstatement{Given an array $\textit{arr}$ of $n$ elements and a list
of query indices, for each query index $i$, count \textbf{how many}
elements to the right of $i$ (not just the nearest one) are strictly
greater than $\textit{arr}[i]$.}

\begin{patternbox}
This problem is grouped alongside the monotonic-stack problems because it
shares their right-to-left scanning direction, but the keyword
\emph{``count how many,''} rather than \emph{``find the nearest,''} is a
signal that a plain monotonic stack is not quite the right tool: a stack
only ever keeps a thin, decreasing/increasing \emph{skyline} of values, not
the full multiset of everything seen so far, so it cannot directly answer
``how many are greater than $X$.'' This calls for a different right-to-left
technique -- a \textbf{Fenwick tree (Binary Indexed Tree) over compressed
ranks} -- that still processes the array in the same right-to-left order
as the rest of this section, but tracks counts rather than a skyline.
\end{patternbox}

\subsection*{Understanding the problem carefully}

Scanning right to left, by the time we reach index $i$, we want to know
how many of the values already seen (i.e.\ to $i$'s right) exceed
$\textit{arr}[i]$. This is a classic ``count of elements greater than $X$
among those inserted so far'' query, repeated $n$ times with the
collection growing by one element after each query -- exactly what a
Fenwick tree, indexed by the \emph{rank} of each value (after coordinate
compression, since values may be arbitrary integers), is built for: it
supports both ``insert a value'' and ``count how many inserted values are
$\le$ some threshold'' in $O(\log n)$ each.

\begin{ideabox}[Fenwick Tree over Compressed Ranks, Scanned Right to Left]
Compress all values in $\textit{arr}$ to ranks $1, \dots, n$ (sorted
order). Scan $\textit{arr}$ from right to left. For each $\textit{arr}[i]$:
query the Fenwick tree for the count of already-inserted values with rank
\emph{greater than} $\textit{arr}[i]$'s rank (equivalently, $n$ minus the
count of values with rank $\le \textit{arr}[i]$'s rank) -- this is exactly
the number of elements to $i$'s right that exceed $\textit{arr}[i]$.
Record this count for index $i$. Then insert $\textit{arr}[i]$'s rank into
the Fenwick tree, so it is correctly counted for any index further to the
left.
\end{ideabox}

\subsection*{Why scanning right to left and inserting afterward is correct}

At the moment we query for index $i$, the Fenwick tree contains exactly
the ranks of elements at indices $i+1, \dots, n-1$ -- because we always
insert $\textit{arr}[i]$ \emph{after} querying for it, and we process
indices in strictly decreasing order, so every index to $i$'s right has
already been inserted, and no index at or to $i$'s left has been inserted
yet. This invariant is what makes the Fenwick tree's count query, at the
moment we use it for index $i$, exactly answer ``how many elements to the
right of $i$ exceed $\textit{arr}[i]$,'' with no extra bookkeeping needed.

\subsection*{Algorithm}

\begin{enumerate}
  \item Compute ranks of all values in $\textit{arr}$ via sorting
        (coordinate compression).
  \item Initialize an empty Fenwick tree of size $n$.
  \item For $i$ from $n-1$ down to $0$:
  \begin{enumerate}
    \item $\textit{countGreater}[i] \gets n - \Call{Query}{\textit{rank}[i]}$
          (count of inserted values with rank $\le \textit{rank}[i]$,
          subtracted from the total inserted so far).
    \item $\Call{Update}{\textit{rank}[i], +1}$ (insert this rank).
  \end{enumerate}
  \item Answer each query index by direct lookup into
        $\textit{countGreater}$.
\end{enumerate}

\subsection*{Dry run}

\dryrun{Take $\textit{arr} = [3,4,2,7,5]$, query index $i=2$ (value
$2$).}

Ranks (1-indexed, smallest $=1$): $3\to2, 4\to3, 2\to1, 7\to5, 5\to4$.

Scanning right to left and inserting:

\begin{itemize}
  \item $i=4$ (value $5$, rank $4$): tree empty, query returns $0$
        inserted so far, so $\textit{countGreater}[4]=0-0=0$. Insert
        rank $4$.
  \item $i=3$ (value $7$, rank $5$): inserted so far $=1$ (rank $4$),
        count with rank $\le5$ is $1$. $\textit{countGreater}[3]=1-1=0$.
        Insert rank $5$.
  \item $i=2$ (value $2$, rank $1$): inserted so far $=2$ (ranks
        $4,5$), count with rank $\le1$ is $0$.
        $\textit{countGreater}[2]=2-0=2$. Insert rank $1$.
\end{itemize}

$\textit{countGreater}[2]=2$: indeed, to the right of index $2$ (value
$2$), the elements are $7$ and $5$, both greater -- $2$ elements,
matching.

\subsection*{C++ implementation}

\begin{lstlisting}
#include <bits/stdc++.h>
using namespace std;

class FenwickTree {
    vector<int> tree;
    int n;
public:
    FenwickTree(int n) : n(n), tree(n + 1, 0) {}

    void update(int i, int delta) {
        for (; i <= n; i += i & (-i)) tree[i] += delta;
    }

    int query(int i) { // count of inserted values with rank <= i
        int sum = 0;
        for (; i > 0; i -= i & (-i)) sum += tree[i];
        return sum;
    }
};

vector<int> countGreaterToRight(vector<int>& arr) {
    int n = arr.size();
    vector<int> sorted = arr;
    sort(sorted.begin(), sorted.end());

    auto rankOf = [&](int val) {
        return lower_bound(sorted.begin(), sorted.end(), val) - sorted.begin() + 1;
    };

    FenwickTree ft(n);
    vector<int> countGreater(n);
    int insertedSoFar = 0;

    for (int i = n - 1; i >= 0; i--) {
        int r = rankOf(arr[i]);
        countGreater[i] = insertedSoFar - ft.query(r);
        ft.update(r, 1);
        insertedSoFar++;
    }
    return countGreater;
}

int main() {
    vector<int> arr = {3, 4, 2, 7, 5};
    auto result = countGreaterToRight(arr);
    for (int x : result) cout << x << " ";
    cout << endl;
    return 0;
}
\end{lstlisting}

\begin{complexitybox}
\textbf{Time Complexity.} Coordinate compression costs $O(n \log n)$. The
main scan performs $n$ Fenwick tree updates and queries, each
$O(\log n)$, giving an overall time complexity of
\[
  O(n \log n).
\]
\textbf{Space Complexity.} $O(n)$ for the Fenwick tree and rank arrays.
\end{complexitybox}

\begin{takeawaybox}
``Count how many'' is a structurally different question from ``find the
nearest,'' even when both scan in the same direction -- a monotonic stack
answers the latter in $O(n)$ total but cannot answer the former at all,
since it discards the very elements a counting query would need. Whenever
a problem asks for a \emph{count} of qualifying elements rather than the
nearest one, reach for a Fenwick tree, segment tree, or balanced
order-statistics structure instead of a plain stack.
\end{takeawaybox}

\section{Trapping Rain Water}
\label{sec:sq-trapping-rain-water}

\problemstatement{We are given $n$ non-negative integers $\textit{height}$
representing an elevation map, where $\textit{height}[i]$ is the height of
a bar of width $1$ at position $i$. Compute how much water is trapped
after raining, assuming water can only be held between bars and runs off
at either open end.}

\begin{patternbox}
\emph{``Water trapped between bars of varying heights''} is a problem
about \textbf{valleys bounded by two walls}: water above any single
position is limited by the \emph{shorter} of the nearest wall to its left
and the nearest wall to its right. Whenever a problem's answer at each
position depends on ``the nearest taller boundary on each side,'' a
monotonic stack -- which is precisely a stack of candidate boundaries,
each taller than the one below it -- is a natural fit, processing the
array once and resolving each ``valley'' as soon as a wall tall enough to
close it off arrives.
\end{patternbox}

\subsection*{Understanding the problem carefully}

Maintain a stack of indices whose heights are non-increasing from bottom
to top (each new index pushed is no taller than the one below it,
otherwise it would immediately start resolving valleys). When we
encounter a bar taller than the index on top of the stack, that means we
have just found a \textbf{right wall} for a valley whose floor is the
popped index, and whose \textbf{left wall} is whatever index remains below
it on the stack. The water trapped above that floor, between the two
walls, is determined by the \emph{shorter} of the two walls (water cannot
rise above the shorter wall without spilling over it) minus the floor's
own height, multiplied by the horizontal distance between the two walls.

\begin{ideabox}[Resolve One Valley Floor per Pop]
Scan left to right. While the current bar is taller than the stack's top:
pop that top index as the \textit{valley floor}. If the stack is now
empty, there is no left wall, so this floor traps no water (stop popping
for this bar). Otherwise, the new stack top is the \textit{left wall}, and
the current index is the \textit{right wall}. The trapped water above the
floor is
\[
  (\min(\textit{height}[\textit{leftWall}],
  \textit{height}[\textit{rightWall}]) - \textit{height}[\textit{floor}])
  \times (\textit{rightWall} - \textit{leftWall} - 1).
\]
Add this to the running total, then continue popping (the same right wall
might close off multiple successively higher floors). Finally, push the
current index.
\end{ideabox}

\subsection*{Why resolving level by level, floor by floor, correctly sums all trapped water}

Every unit of trapped water sits directly above some specific bar (the
``floor'') and below the line connecting the nearest sufficiently tall
walls on either side. By processing each floor exactly once -- the moment
a right wall taller than it is found, with whatever remains on the stack
serving as the left wall -- we account for the entire water column above
that floor in one $O(1)$ computation. Because indices are only ever pushed
once and popped at most once, every floor's contribution is computed
exactly once, and the total accumulated is exactly the total trapped
water, with no double-counting and no gaps.

\subsection*{Algorithm}

\begin{enumerate}
  \item Initialize an empty stack and $\textit{totalWater} \gets 0$.
  \item For $i$ from $0$ to $n-1$:
  \begin{enumerate}
    \item While the stack is not empty and $\textit{height}[i] >
          \textit{height}[\text{stack's top}]$:
    \begin{enumerate}
      \item Pop the top as $\textit{floor}$.
      \item If the stack is now empty: break (no left wall).
      \item $\textit{leftWall} \gets$ the new top; $\textit{width}
            \gets i - \textit{leftWall} - 1$;
            $\textit{boundedHeight} \gets
            \min(\textit{height}[\textit{leftWall}],
            \textit{height}[i]) - \textit{height}[\textit{floor}]$.
      \item $\textit{totalWater} \gets \textit{totalWater} +
            \textit{width} \times \textit{boundedHeight}$.
    \end{enumerate}
    \item Push $i$.
  \end{enumerate}
  \item Return \textit{totalWater}.
\end{enumerate}

\subsection*{Dry run}

\dryrun{Take $\textit{height} = [3,0,2,0,4]$.}

\begin{itemize}
  \item $i=0$ ($h=3$): stack empty. Push $0$. Stack $=[0]$.
  \item $i=1$ ($h=0$): $0 \le \textit{height}[0]=3$, no pop. Push $1$.
        Stack $=[0,1]$.
  \item $i=2$ ($h=2$): $2 > \textit{height}[1]=0$, pop floor $=1$. Stack
        $=[0]$, not empty: leftWall$=0$. Width $=2-0-1=1$.
        boundedHeight $=\min(3,2)-0=2$. Water $+=1\times2=2$.
        $\textit{totalWater}=2$. Now $2 \le \textit{height}[0]=3$, stop
        popping. Push $2$. Stack $=[0,2]$.
  \item $i=3$ ($h=0$): $0\le \textit{height}[2]=2$, no pop. Push $3$.
        Stack $=[0,2,3]$.
  \item $i=4$ ($h=4$): $4>\textit{height}[3]=0$, pop floor$=3$. Left
        wall $=2$. Width $=4-2-1=1$. boundedHeight
        $=\min(2,4)-0=2$. Water $+=1\times2=2$. Total$=4$. Continue:
        $4>\textit{height}[2]=2$, pop floor$=2$. Left wall$=0$. Width
        $=4-0-1=3$. boundedHeight $=\min(3,4)-2=1$. Water
        $+=3\times1=3$. Total$=7$. Now $4>\textit{height}[0]=3$, pop
        floor$=0$. Stack empty -- break (no left wall). Push $4$.
\end{itemize}

Total trapped water: $7$.

\subsection*{C++ implementation}

\begin{lstlisting}
#include <bits/stdc++.h>
using namespace std;

int trapRainWater(vector<int>& height) {
    int n = height.size();
    stack<int> st;
    long long totalWater = 0;

    for (int i = 0; i < n; i++) {
        while (!st.empty() && height[i] > height[st.top()]) {
            int floorIdx = st.top();
            st.pop();
            if (st.empty()) break;

            int leftWall = st.top();
            int width = i - leftWall - 1;
            int boundedHeight = min(height[leftWall], height[i]) - height[floorIdx];
            totalWater += (long long)width * boundedHeight;
        }
        st.push(i);
    }
    return totalWater;
}

int main() {
    vector<int> height = {3, 0, 2, 0, 4};
    cout << "Trapped water: " << trapRainWater(height) << endl;
    return 0;
}
\end{lstlisting}

\begin{complexitybox}
\textbf{Time Complexity.} Each index is pushed once and popped at most
once, giving
\[
  O(n).
\]
\textbf{Space Complexity.} $O(n)$ for the stack in the worst case. (A
two-pointer technique solves this same problem in $O(1)$ extra space, by
tracking running left-max and right-max boundaries directly -- worth
knowing as a further optimization, though the stack-based derivation here
is the one that generalizes most directly to the histogram-area problems
in Sections~\ref{sec:sq-largest-rectangle-histogram}--\ref{sec:sq-maximal-rectangle}.)
\end{complexitybox}

\begin{takeawaybox}
``Water/area bounded by the nearest sufficiently tall walls on both
sides'' is solved by a monotonic stack that resolves one bounded region
per pop, the instant a wall tall enough to close it off arrives. This same
skeleton -- pop a floor, look at what remains as the left boundary, the
trigger as the right boundary, compute a bounded quantity, accumulate --
reappears in Largest Rectangle in Histogram
(Section~\ref{sec:sq-largest-rectangle-histogram}).
\end{takeawaybox}

\section{Sum of Subarray Minimums}
\label{sec:sq-sum-subarray-minimums}

\problemstatement{Given an array $\textit{arr}$ of $n$ integers, consider
every one of its $\binom{n+1}{2}$ contiguous subarrays. Compute the sum,
over all of them, of each subarray's minimum element (the result can be
large, so it is typically returned modulo $10^9+7$).}

\begin{patternbox}
\emph{``Sum, over all subarrays, of each subarray's minimum (or
maximum)''} is a strong signal to flip the entire computation around:
instead of asking, for each subarray, ``what is its minimum,'' ask, for
each \textbf{element}, ``in how many subarrays is \emph{this} element the
minimum.'' This contribution-counting reformulation, combined with a
monotonic stack to compute exactly that count for every element in
$O(n)$ total, avoids ever enumerating the $O(n^2)$ subarrays directly.
\end{patternbox}

\subsection*{Understanding the problem carefully}

For element $\textit{arr}[i]$ to be the minimum of some subarray, that
subarray must be entirely contained within the maximal range around $i$
over which $\textit{arr}[i]$ remains the smallest value -- bounded on the
left by the nearest strictly smaller element, and on the right by the
nearest smaller-or-equal element (using a strict boundary on one side and
a non-strict boundary on the other is essential to avoid double-counting
subarrays when duplicate values are present, by consistently attributing
each subarray's minimum to a single, well-defined occurrence among any
tied values).

\begin{ideabox}[{Count Subarrays Where $\textit{arr}[i]$ Is the Minimum}]
For each $i$, let $\textit{left}[i]$ be the number of valid choices for a
subarray's left endpoint such that every element from that endpoint up to
$i$ is $\ge \textit{arr}[i]$ -- this equals $i - \textit{PLE}(i)$, where
$\textit{PLE}(i)$ is the index of the nearest element to the left that is
\emph{strictly less} than $\textit{arr}[i]$ (or $-1$ if none). Let
$\textit{right}[i]$ be the number of valid choices for the right endpoint
such that every element from $i$ to that endpoint is $\ge
\textit{arr}[i]$ -- this equals $\textit{NLE}(i) - i$, where
$\textit{NLE}(i)$ is the nearest element to the right that is \emph{less
than or equal to} $\textit{arr}[i]$ (or $n$ if none). Then
$\textit{arr}[i]$ is the minimum of exactly $\textit{left}[i] \times
\textit{right}[i]$ subarrays (any combination of a valid left endpoint and
a valid right endpoint).
\end{ideabox}

\subsection*{Why the strict/non-strict boundary asymmetry avoids double-counting}

If both boundaries used the same (either both strict, or both
non-strict) comparison, a subarray containing two equal minimum values
would be counted once for \emph{each} occurrence of that value, inflating
the sum. Using strictly-less on the left and less-or-equal on the right
ensures that, among several equal candidate minimum positions within a
subarray, only the \textbf{leftmost} one ``claims'' that subarray as
contributing to its count: the leftmost equal occurrence's right boundary
search will stop at the very next equal (or smaller) value, while any
later equal occurrence's left boundary search will pass straight over
the earlier equal value (since it is not \emph{strictly} less) and only
stop at a genuinely smaller one further left -- so the two occurrences
never both count the same subarray.

\subsection*{Algorithm}

\begin{enumerate}
  \item Compute $\textit{left}[i]$ for every $i$ via a monotonic
        increasing stack scanned left to right: pop while the stack's top
        value $\ge \textit{arr}[i]$; $\textit{left}[i] \gets i - (\text{
        stack empty} ? -1 : \text{top index})$; push $i$.
  \item Compute $\textit{right}[i]$ for every $i$ via a monotonic stack
        scanned right to left: pop while the stack's top value $>
        \textit{arr}[i]$; $\textit{right}[i] \gets (\text{stack empty} ?
        n : \text{top index}) - i$; push $i$.
  \item Sum $\textit{arr}[i] \times \textit{left}[i] \times
        \textit{right}[i]$ over all $i$, modulo $10^9+7$.
\end{enumerate}

\subsection*{Dry run}

\dryrun{Take $\textit{arr} = [3,1,2,4]$.}

\textbf{Left} (pop while top $\ge \textit{arr}[i]$): $\textit{left} =
[1,2,1,1]$ (computed via the standard PLE scan, e.g.\ at $i=2$, value
$2$: the stack's top is index $1$, value $1 < 2$, so no pop;
$\textit{left}[2]=2-1=1$).

\textbf{Right} (pop while top $> \textit{arr}[i]$): $\textit{right} =
[1,3,2,1]$ (e.g.\ at $i=1$, value $1$: scanning right to left, both
index $2$ (value $2$) and index $3$... after appropriate pops the stack
empties, giving $\textit{right}[1]=4-1=3$).

Contributions: $3 \times 1 \times 1=3$; $1\times2\times3=6$;
$2\times1\times2=4$; $4\times1\times1=4$. Total $=3+6+4+4=17$.

\subsection*{C++ implementation}

\begin{lstlisting}
#include <bits/stdc++.h>
using namespace std;
const int MOD = 1e9 + 7;

int sumSubarrayMins(vector<int>& arr) {
    int n = arr.size();
    vector<int> left(n), right(n);
    stack<int> st;

    for (int i = 0; i < n; i++) {
        while (!st.empty() && arr[st.top()] >= arr[i]) st.pop();
        left[i] = i - (st.empty() ? -1 : st.top());
        st.push(i);
    }

    while (!st.empty()) st.pop();

    for (int i = n - 1; i >= 0; i--) {
        while (!st.empty() && arr[st.top()] > arr[i]) st.pop();
        right[i] = (st.empty() ? n : st.top()) - i;
        st.push(i);
    }

    long long total = 0;
    for (int i = 0; i < n; i++) {
        total = (total + (long long)arr[i] * left[i] % MOD * right[i]) % MOD;
    }
    return total;
}

int main() {
    vector<int> arr = {3, 1, 2, 4};
    cout << "Sum of subarray minimums: " << sumSubarrayMins(arr) << endl;
    return 0;
}
\end{lstlisting}

\begin{complexitybox}
\textbf{Time Complexity.} Two monotonic-stack passes, each $O(n)$, giving
\[
  O(n).
\]
\textbf{Space Complexity.} $O(n)$ for the stacks and the
\textit{left}/\textit{right} arrays.
\end{complexitybox}

\begin{takeawaybox}
``Sum/count over all subarrays of [some per-subarray extremum]'' is
solved not by enumerating subarrays, but by computing, for each element,
\emph{how many} subarrays it is the extremum of -- a count derived from
the nearest strictly-more-extreme boundaries on each side, found with two
monotonic-stack passes. This contribution-counting reformulation is one
of the most powerful general techniques to recognize in this chapter.
\end{takeawaybox}

\section{Sum of Subarray Ranges}
\label{sec:sq-sum-subarray-ranges}

\problemstatement{Given an array $\textit{arr}$, the \textbf{range} of a
subarray is defined as its maximum element minus its minimum element.
Compute the sum of ranges over all contiguous subarrays.}

\begin{patternbox}
A subarray's range is, by definition, $\max - \min$. Summing ranges over
all subarrays is therefore exactly
\[
  \left(\sum_{\text{subarrays}} \max\right) -
  \left(\sum_{\text{subarrays}} \min\right),
\]
and the second term is precisely Section~\ref{sec:sq-sum-subarray-minimums}'s
problem. The first term is its exact mirror image -- ``sum of subarray
maximums'' -- solved by the same contribution-counting technique with
every comparison flipped. No new derivation is needed; this problem is a
direct combination of two near-identical applications of the same idea.
\end{patternbox}

\subsection*{Understanding the problem carefully}

Linearity of summation lets us split the sum of $(\max - \min)$ over all
subarrays into a sum of $\max$ over all subarrays, minus a sum of
$\min$ over all subarrays, computed completely independently. Each piece
is solved by the contribution-counting technique from
Section~\ref{sec:sq-sum-subarray-minimums}: for the minimum part, exactly
as before; for the maximum part, the same left/right boundary counting,
but searching for the nearest strictly-\emph{greater} element (instead of
strictly-less) on one side and greater-or-equal on the other.

\begin{ideabox}[Reuse Twice, with Comparisons Flipped for the Maximum Half]
Compute $\textit{sumOfMins}$ exactly as in
Algorithm~\ref{sec:sq-sum-subarray-minimums}. Compute
$\textit{sumOfMaxs}$ with the same algorithm, but with stack comparisons
flipped: the left pass pops while the stack's top value $\le
\textit{arr}[i]$ (finding the nearest strictly \emph{greater} element to
the left), and the right pass pops while the stack's top value $<
\textit{arr}[i]$ (finding the nearest greater-or-equal element to the
right). Return $\textit{sumOfMaxs} - \textit{sumOfMins}$.
\end{ideabox}

\subsection*{Why flipping every comparison correctly computes the maximum-side sum}

This is a direct restatement of the next-greater/next-smaller duality
established in Section~\ref{sec:sq-next-smaller}: ``how many subarrays is
$\textit{arr}[i]$ the maximum of'' is structurally identical to ``how many
subarrays is $\textit{arr}[i]$ the minimum of,'' with every $<$ replaced
by $>$ and vice versa. The strict/non-strict asymmetry between the two
boundary searches is still required, for the same reason as
Section~\ref{sec:sq-sum-subarray-minimums}: to avoid double-counting
subarrays containing duplicate maximum values.

\subsection*{Algorithm}

\begin{enumerate}
  \item Compute $\textit{sumOfMins}$ using
        Algorithm~\ref{sec:sq-sum-subarray-minimums} unchanged.
  \item Compute $\textit{sumOfMaxs}$ using the same algorithm structure,
        with left-pass pops on $\le$ and right-pass pops on $<$ (instead
        of $\ge$ and $>$ respectively).
  \item Return $\textit{sumOfMaxs} - \textit{sumOfMins}$.
\end{enumerate}

\subsection*{Dry run}

\dryrun{Take $\textit{arr} = [1,2,3]$.}

$\textit{sumOfMins}$: subarrays and their minimums: $[1]{=}1$, $[2]{=}2$,
$[3]{=}3$, $[1,2]{=}1$, $[2,3]{=}2$, $[1,2,3]{=}1$. Sum $=1+2+3+1+2+1=10$.

$\textit{sumOfMaxs}$: maximums: $1,2,3,2,3,3$. Sum
$=1+2+3+2+3+3=14$.

Sum of ranges $=14-10=4$. Checking directly: ranges are $0,0,0,1,1,2$,
summing to $4$. Matches.

\subsection*{C++ implementation}

\begin{lstlisting}
#include <bits/stdc++.h>
using namespace std;

long long sumOfMins(vector<int>& arr) {
    int n = arr.size();
    vector<int> left(n), right(n);
    stack<int> st;

    for (int i = 0; i < n; i++) {
        while (!st.empty() && arr[st.top()] >= arr[i]) st.pop();
        left[i] = i - (st.empty() ? -1 : st.top());
        st.push(i);
    }
    while (!st.empty()) st.pop();
    for (int i = n - 1; i >= 0; i--) {
        while (!st.empty() && arr[st.top()] > arr[i]) st.pop();
        right[i] = (st.empty() ? n : st.top()) - i;
        st.push(i);
    }

    long long total = 0;
    for (int i = 0; i < n; i++) total += (long long)arr[i] * left[i] * right[i];
    return total;
}

long long sumOfMaxs(vector<int>& arr) {
    int n = arr.size();
    vector<int> left(n), right(n);
    stack<int> st;

    for (int i = 0; i < n; i++) {
        while (!st.empty() && arr[st.top()] <= arr[i]) st.pop();
        left[i] = i - (st.empty() ? -1 : st.top());
        st.push(i);
    }
    while (!st.empty()) st.pop();
    for (int i = n - 1; i >= 0; i--) {
        while (!st.empty() && arr[st.top()] < arr[i]) st.pop();
        right[i] = (st.empty() ? n : st.top()) - i;
        st.push(i);
    }

    long long total = 0;
    for (int i = 0; i < n; i++) total += (long long)arr[i] * left[i] * right[i];
    return total;
}

long long sumSubarrayRanges(vector<int>& arr) {
    return sumOfMaxs(arr) - sumOfMins(arr);
}

int main() {
    vector<int> arr = {1, 2, 3};
    cout << "Sum of subarray ranges: " << sumSubarrayRanges(arr) << endl;
    return 0;
}
\end{lstlisting}

\begin{complexitybox}
\textbf{Time Complexity.} Four monotonic-stack passes total (two for
minimums, two for maximums), each $O(n)$, giving
\[
  O(n).
\]
\textbf{Space Complexity.} $O(n)$.
\end{complexitybox}

\begin{takeawaybox}
When a quantity decomposes linearly (here, range $= \max - \min$),
compute each piece with whatever technique already solves it, and combine
at the end -- there is no need to find a single unified algorithm that
handles $\max$ and $\min$ simultaneously when summing them separately and
subtracting is just as correct and far simpler to reason about.
\end{takeawaybox}

\section{Asteroid Collision}
\label{sec:sq-asteroid-collision}

\problemstatement{We are given an array $\textit{asteroids}$, where each
element's absolute value is an asteroid's size and its sign is its
direction of travel (positive moves right, negative moves left; all move
at the same speed). When two asteroids meet, the smaller one explodes; if
both are the same size, both explode. Asteroids moving in the same
direction never meet. Return the state of the asteroids after all
collisions have finished.}

\begin{patternbox}
\emph{``Process a sequence left to right, where a new element can
repeatedly `defeat' and remove previously-kept elements (or be defeated
by them) based on a comparison''} is, structurally, the same skeleton as
next greater element (Section~\ref{sec:sq-next-greater-1}): a stack of
``currently surviving'' items, where each new arrival potentially knocks
out some suffix of what's already on the stack before either joining it
or being knocked out itself.
\end{patternbox}

\subsection*{Understanding the problem carefully}

A collision can only happen between a \textbf{right-moving} asteroid that
was placed earlier and a \textbf{left-moving} asteroid that arrives later
(two right-movers never catch up to each other since they move at the
same speed; a left-mover followed by a right-mover are moving apart and
never meet). So, scanning left to right and keeping a stack of asteroids
that have survived so far: whenever a new \emph{left-moving} asteroid
arrives and the stack's top is a surviving \emph{right-moving} asteroid,
they are on a collision course, and we must resolve who survives --
possibly triggering a chain of further collisions if the new asteroid is
large enough to destroy multiple right-movers in a row.

\begin{ideabox}[Stack of Surviving Asteroids, with Chain Collisions]
Maintain a stack. For a right-moving asteroid (positive), simply push it
-- no immediate collision is possible yet. For a left-moving asteroid
(negative): while the stack's top is positive (a right-mover) and smaller
in size than the incoming asteroid, pop it (it loses). If the stack's
top is positive and exactly equal in size, pop it and the incoming
asteroid is also destroyed (mutual annihilation) -- stop processing this
asteroid. If the stack's top is positive and larger, the incoming asteroid
is destroyed -- stop processing this asteroid, do not push it. If the
stack becomes empty or its top is negative (no more right-movers to
collide with), and the incoming asteroid was never destroyed, push it.
\end{ideabox}

\subsection*{Why processing collisions as a chain (not just one comparison) is correct}

A single incoming left-moving asteroid might be large enough to destroy
\emph{several} consecutive right-moving asteroids already on the stack
before either finally being destroyed itself or surviving to be pushed.
Each such destruction is resolved by exactly the comparison rule above,
applied repeatedly -- this is precisely the same ``pop everything the new
arrival invalidates'' loop structure as a monotonic stack, except the
``invalidation'' rule here is a size comparison between asteroids of
opposite direction rather than a simple value comparison, and the loop
must also account for the three-way outcome (incoming destroyed, top
destroyed, or both destroyed) rather than a single binary pop/no-pop
decision.

\subsection*{Algorithm}

\begin{enumerate}
  \item Initialize an empty stack.
  \item For each asteroid $a$ in $\textit{asteroids}$:
  \begin{enumerate}
    \item $\textit{alive} \gets \text{true}$.
    \item While $\textit{alive}$, $a < 0$, the stack is not empty, and
          the stack's top is $>0$:
    \begin{itemize}
      \item If top $< -a$: pop (the top asteroid is destroyed); continue
            the loop.
      \item Else if top $= -a$: pop; $\textit{alive} \gets \text{false}$
            (mutual destruction).
      \item Else (top $> -a$): $\textit{alive} \gets \text{false}$
            (the incoming asteroid is destroyed).
    \end{itemize}
    \item If $\textit{alive}$: push $a$.
  \end{enumerate}
  \item Return the stack's contents, bottom to top.
\end{enumerate}

\subsection*{Dry run}

\dryrun{Take $\textit{asteroids} = [5, 10, -5]$.}

\begin{itemize}
  \item $a=5$: positive, push. Stack $=[5]$.
  \item $a=10$: positive, push. Stack $=[5,10]$.
  \item $a=-5$: negative. Top is $10>0$. Is $10 < 5$? No. Is $10=5$?
        No. So top ($10$) $> 5$: the incoming asteroid ($-5$) is
        destroyed. \textit{alive} $=$ false. Do not push.
\end{itemize}

Final stack: $[5,10]$.

Now take $\textit{asteroids} = [8,-8]$: $a=8$ pushed. $a=-8$: top is $8$,
and $8 = 8$: pop, mutual destruction, \textit{alive}$=$false. Final
stack: $[]$ (both annihilated).

\subsection*{C++ implementation}

\begin{lstlisting}
#include <bits/stdc++.h>
using namespace std;

vector<int> asteroidCollision(vector<int>& asteroids) {
    vector<int> st; // used as a stack via back()/push_back()/pop_back()

    for (int a : asteroids) {
        bool alive = true;
        while (alive && a < 0 && !st.empty() && st.back() > 0) {
            if (st.back() < -a) {
                st.pop_back(); // top destroyed, continue checking
            } else if (st.back() == -a) {
                st.pop_back();
                alive = false; // mutual destruction
            } else {
                alive = false; // incoming asteroid destroyed
            }
        }
        if (alive) st.push_back(a);
    }
    return st;
}

int main() {
    vector<int> asteroids = {5, 10, -5};
    for (int x : asteroidCollision(asteroids)) cout << x << " ";
    cout << endl;
    return 0;
}
\end{lstlisting}

\begin{complexitybox}
\textbf{Time Complexity.} Each asteroid is pushed at most once and popped
at most once across the entire scan, giving
\[
  O(n).
\]
\textbf{Space Complexity.} $O(n)$ for the stack in the worst case (no
collisions occur).
\end{complexitybox}

\begin{takeawaybox}
Monotonic-stack-style ``pop everything the new arrival invalidates''
logic generalizes beyond simple value comparisons: here the invalidation
rule is a three-way size comparison between opposite-direction asteroids,
but the overall shape -- a while-loop popping a suffix of the stack, then
conditionally pushing the new arrival -- is unchanged from
Section~\ref{sec:sq-next-greater-1}.
\end{takeawaybox}

\section{Remove K Digits}
\label{sec:sq-remove-k-digits}

\problemstatement{Given a non-negative integer represented as a string
$\textit{num}$, and an integer $k$, remove exactly $k$ digits from
$\textit{num}$ so that the remaining digits, kept in their original
relative order, form the \textbf{smallest possible} number. The result
must not have leading zeros (unless it is exactly \texttt{"0"}).}

\begin{patternbox}
\emph{``Remove $k$ elements (keeping relative order) to make the result as
small/large as possible''} is a greedy-with-a-stack pattern: a number is
smaller when an earlier (more significant) digit is smaller, so whenever
we see a digit smaller than what we have just placed, removing that
just-placed (now regretted) larger digit is always an improvement -- a
monotonic increasing stack of kept digits, with removals funded by the
budget $k$, captures exactly this.
\end{patternbox}

\subsection*{Understanding the problem carefully}

Among any two digit strings of the same length, the one with the smaller
digit at the first position where they differ is the smaller number --
so, greedily, whenever we have the opportunity to remove a digit and
thereby make some earlier-kept digit smaller (by exposing a smaller digit
to take its place at that position), we should take it, as long as we
still have removals left in our budget $k$.

\begin{ideabox}[Monotonic Increasing Stack of Kept Digits, Funded by Budget $k$]
Scan $\textit{num}$ left to right, maintaining a stack of digits kept so
far. For each new digit $d$: while $k>0$ and the stack is non-empty and
its top digit exceeds $d$, pop the top (spend one unit of $k$) -- this
digit was a mistake to keep, now that a smaller one is available to take
its place. Then push $d$. After the scan, if $k$ still has unused budget,
remove that many digits from the \emph{end} of the stack (the largest
remaining suffix, since no smaller digit ever showed up to justify
removing them earlier). Finally, strip any leading zeros, and return
\texttt{"0"} if the result is empty.
\end{ideabox}

\subsection*{Why popping a larger top whenever a smaller digit arrives is optimal}

Removing a digit only ever helps if it lets an earlier, larger digit be
replaced (in significance/position) by a smaller later digit -- and the
earliest opportunity to do this should always be taken, since a smaller
digit at a \emph{more significant} (more leftward) position always
beats any improvement achievable further right. This is exactly the same
exchange-argument logic as every monotonic-stack pop rule in this
chapter: keeping a larger digit when a smaller one is available and
removals remain in budget can never be better than removing it, since the
resulting number would be lexicographically (and numerically, for
same-length strings) larger or equal in every case.

\subsection*{Algorithm}

\begin{enumerate}
  \item Initialize an empty stack of digits.
  \item For each digit $d$ in $\textit{num}$:
  \begin{enumerate}
    \item While $k>0$ and the stack is non-empty and its top $>d$: pop;
          $k \gets k-1$.
    \item Push $d$.
  \end{enumerate}
  \item While $k>0$: pop from the end of the stack; $k \gets k-1$.
  \item Strip leading zeros from the stack's contents (read bottom to
        top). If the result is empty, return \texttt{"0"}.
\end{enumerate}

\subsection*{Dry run}

\dryrun{Take $\textit{num} = \texttt{"1432219"}$, $k=3$.}

\begin{itemize}
  \item \texttt{1}: stack empty, push. Stack $=[1]$.
  \item \texttt{4}: top $1<4$, no pop. Push. Stack $=[1,4]$.
  \item \texttt{3}: top $4>3$ and $k=3>0$: pop $4$, $k=2$. New top
        $1<3$, stop. Push $3$. Stack $=[1,3]$.
  \item \texttt{2}: top $3>2$ and $k=2>0$: pop $3$, $k=1$. New top
        $1<2$, stop. Push $2$. Stack $=[1,2]$.
  \item \texttt{2}: top $2 \not> 2$, no pop. Push. Stack $=[1,2,2]$.
  \item \texttt{1}: top $2>1$ and $k=1>0$: pop $2$, $k=0$. New top
        $2>1$ but $k=0$ now, stop. Push $1$. Stack $=[1,2,1]$.
  \item \texttt{9}: $k=0$, no pop. Push. Stack $=[1,2,1,9]$.
\end{itemize}

$k=0$, no trailing removal needed. Result (no leading zeros to strip):
\texttt{"1219"}.

\subsection*{C++ implementation}

\begin{lstlisting}
#include <bits/stdc++.h>
using namespace std;

string removeKdigits(string num, int k) {
    string st; // used as a stack of characters

    for (char d : num) {
        while (k > 0 && !st.empty() && st.back() > d) {
            st.pop_back();
            k--;
        }
        st.push_back(d);
    }

    while (k > 0 && !st.empty()) {
        st.pop_back();
        k--;
    }

    int i = 0;
    while (i < (int)st.size() - 1 && st[i] == '0') i++;
    st = st.substr(i);

    return st.empty() ? "0" : st;
}

int main() {
    cout << removeKdigits("1432219", 3) << endl;
    return 0;
}
\end{lstlisting}

\begin{complexitybox}
\textbf{Time Complexity.} Each digit is pushed once and popped at most
once, giving
\[
  O(n).
\]
\textbf{Space Complexity.} $O(n)$ for the stack.
\end{complexitybox}

\begin{takeawaybox}
``Remove exactly $k$ elements to minimize/maximize the resulting
sequence, preserving relative order'' is solved by a monotonic stack
where the popping budget is explicitly $k$ -- once $k$ reaches zero, the
stack stops discarding even objectively ``bad'' elements, and any
leftover budget at the end is spent removing from the tail, since nothing
better ever arrived to justify earlier removals there.
\end{takeawaybox}

\section{Largest Rectangle in Histogram}
\label{sec:sq-largest-rectangle-histogram}

\problemstatement{Given an array $\textit{heights}$ representing the
heights of adjacent bars of width $1$ in a histogram, find the area of
the \textbf{largest rectangle} that can be formed using contiguous bars
(the rectangle's height is limited by the shortest bar it spans).}

\begin{patternbox}
This is the contribution-counting idea from
Section~\ref{sec:sq-sum-subarray-minimums}, repurposed to
\textbf{maximize} rather than sum: for every bar $i$, the largest
rectangle that uses $\textit{heights}[i]$ as its limiting height spans
exactly the maximal contiguous range around $i$ over which
$\textit{heights}[i]$ is the minimum -- found by the exact same nearest
strictly-smaller-element boundaries on each side.
\end{patternbox}

\subsection*{Understanding the problem carefully}

For a rectangle of height $\textit{heights}[i]$ to be valid, every bar it
spans must be at least $\textit{heights}[i]$ tall (otherwise the
rectangle would poke above some shorter bar within its width). So the
widest rectangle using $\textit{heights}[i]$ as its height extends left
until hitting a strictly shorter bar (or the histogram's edge), and right
until hitting a strictly shorter bar (or the edge) -- exactly the
$\textit{left}[i]$ and $\textit{right}[i]$ boundaries from
Section~\ref{sec:sq-sum-subarray-minimums}, just used to compute one
candidate area per bar and take the maximum, rather than summing
contributions.

\begin{ideabox}[One Pop, One Area, via a Single Monotonic Pass]
Maintain a stack of indices with strictly increasing heights. Scan left
to right (appending a sentinel height of $0$ at the very end to force
every remaining bar to be resolved). When the current height is less than
the stack's top height, pop that top index $j$: the rectangle with height
$\textit{heights}[j]$ has width $i - (\text{new top index}) - 1$ (or
simply $i$ if the stack becomes empty, meaning $j$'s rectangle could
extend all the way to the histogram's left edge). Compute this area and
update the running maximum. Continue popping until the current height no
longer triggers a pop, then push the current index.
\end{ideabox}

\subsection*{Why a single left-to-right pass (with a sentinel) suffices}

This combines the boundary-finding idea from
Section~\ref{sec:sq-sum-subarray-minimums} with the floor-resolution idea
from Trapping Rain Water (Section~\ref{sec:sq-trapping-rain-water}): each
bar's rectangle is fully resolved -- both its left and right boundary
determined -- in a single pop, because by the time bar $j$ is popped, the
current index $i$ is exactly its right boundary (the nearest strictly
shorter bar to the right), and the stack's new top after the pop is
exactly its left boundary (the nearest still-standing, and therefore
strictly shorter or absent, bar to the left). The sentinel height $0$ at
the end guarantees every bar still on the stack when the scan finishes is
popped and resolved, rather than left unresolved.

\subsection*{Algorithm}

\begin{enumerate}
  \item Append a sentinel value $0$ to the end of $\textit{heights}$
        (conceptually); initialize an empty stack and
        $\textit{maxArea} \gets 0$.
  \item For $i$ from $0$ to $n$ (inclusive of the sentinel position):
  \begin{enumerate}
    \item While the stack is not empty and
          $\textit{heights}[i] < \textit{heights}[\text{stack's top}]$:
    \begin{itemize}
      \item Pop the top as $j$.
      \item $\textit{width} \gets$ (stack empty $? i :$
            $i - \text{new top} - 1$).
      \item $\textit{maxArea} \gets \max(\textit{maxArea},
            \textit{heights}[j] \times \textit{width})$.
    \end{itemize}
    \item Push $i$.
  \end{enumerate}
  \item Return \textit{maxArea}.
\end{enumerate}

\subsection*{Dry run}

\dryrun{Take $\textit{heights} = [2,1,5,6,2,3]$ (sentinel $0$ appended at
index $6$).}

\begin{itemize}
  \item $i=0$ ($h=2$): stack empty, push. Stack $=[0]$.
  \item $i=1$ ($h=1$): $1<2$, pop $0$. Stack empty: width $=1$. Area
        $=2\times1=2$. $\textit{maxArea}=2$. Push $1$. Stack $=[1]$.
  \item $i=2$ ($h=5$): $5 \not< 1$, no pop. Push. Stack $=[1,2]$.
  \item $i=3$ ($h=6$): $6 \not< 5$, no pop. Push. Stack $=[1,2,3]$.
  \item $i=4$ ($h=2$): $2<6$, pop $3$: width $=4-2-1=1$. Area
        $=6\times1=6$. $\textit{maxArea}=6$. $2<5$, pop $2$: width
        $=4-1-1=2$. Area $=5\times2=10$. $\textit{maxArea}=10$. $2
        \not< 1$, stop. Push $4$. Stack $=[1,4]$.
  \item $i=5$ ($h=3$): $3 \not< 2$, no pop. Push. Stack $=[1,4,5]$.
  \item $i=6$ ($h=0$, sentinel): $0<3$, pop $5$: width $=6-4-1=1$. Area
        $=3$. $0<2$, pop $4$: width $=6-1-1=4$. Area $=8$. $0<1$, pop
        $1$: stack empty: width $=6$. Area $=6$. None exceed $10$.
\end{itemize}

The largest rectangle has area $\boxed{10}$ (bars of height $5$ and $6$
at indices $2,3$, using height $5$ across width $2$).

\subsection*{C++ implementation}

\begin{lstlisting}
#include <bits/stdc++.h>
using namespace std;

int largestRectangleArea(vector<int>& heights) {
    heights.push_back(0); // sentinel
    int n = heights.size();
    stack<int> st;
    long long maxArea = 0;

    for (int i = 0; i < n; i++) {
        while (!st.empty() && heights[i] < heights[st.top()]) {
            int j = st.top();
            st.pop();
            int width = st.empty() ? i : i - st.top() - 1;
            maxArea = max(maxArea, (long long)heights[j] * width);
        }
        st.push(i);
    }
    heights.pop_back(); // restore original array
    return maxArea;
}

int main() {
    vector<int> heights = {2, 1, 5, 6, 2, 3};
    cout << "Largest rectangle area: "
         << largestRectangleArea(heights) << endl;
    return 0;
}
\end{lstlisting}

\begin{complexitybox}
\textbf{Time Complexity.} Each index is pushed once and popped at most
once, giving
\[
  O(n).
\]
\textbf{Space Complexity.} $O(n)$ for the stack.
\end{complexitybox}

\begin{takeawaybox}
``Largest rectangle/area limited by the shortest bar within some span''
is solved by the same boundary-finding idea as
Section~\ref{sec:sq-sum-subarray-minimums}, condensed into a single
left-to-right pass with a sentinel, since we only need each bar's
\emph{maximum} possible area, not a sum over every possible sub-span.
This single-pass technique is reused directly in Maximal Rectangle
(Section~\ref{sec:sq-maximal-rectangle}) to extend from a 1D histogram to
a full 2D binary matrix.
\end{takeawaybox}

\section{Maximal Rectangle}
\label{sec:sq-maximal-rectangle}

\problemstatement{Given an $n \times m$ binary matrix, find the area of
the largest rectangle that contains only $1$s.}

\begin{patternbox}
\emph{``Largest all-$1$s rectangle in a binary matrix''} is Largest
Rectangle in Histogram (Section~\ref{sec:sq-largest-rectangle-histogram})
applied \textbf{once per row}, where each row defines a histogram of
``how tall a column of consecutive $1$s ends here.'' The keyword to
notice: a 2D ``maximal rectangle of some property'' problem often reduces
to a 1D technique applied row by row, with state carried forward between
rows.
\end{patternbox}

\subsection*{Understanding the problem carefully}

For each row $r$, define $\textit{height}[j]$ as the number of
consecutive $1$s ending at row $r$ in column $j$, looking upward (so
$\textit{height}[j] = 0$ if $\textit{matrix}[r][j] = 0$, and
$\textit{height}[j] = \textit{height}_{\text{previous row}}[j] + 1$ if
$\textit{matrix}[r][j] = 1$). This $\textit{height}$ array is exactly a
histogram: a rectangle of all $1$s that has its \emph{bottom edge} on row
$r$ corresponds exactly to a rectangle in this histogram, since the
histogram's height at column $j$ tells us exactly how far upward a solid
column of $1$s extends from row $r$.

\begin{ideabox}[Run Largest-Rectangle-in-Histogram Once per Row]
Maintain a running $\textit{height}$ array across rows, updating it as
described above after moving to each new row. After updating
$\textit{height}$ for row $r$, run
Algorithm~\ref{sec:sq-largest-rectangle-histogram}'s histogram algorithm
on it, and record the resulting maximum area as a candidate for the
overall answer. After processing every row, the largest area seen across
all rows is the answer.
\end{ideabox}

\subsection*{Why checking every row's histogram captures every possible rectangle}

Every all-$1$s rectangle in the matrix has \emph{some} specific bottom
row $r$. When we process row $r$, the running $\textit{height}$ array at
that point correctly reflects, for every column, how far a solid run of
$1$s extends upward from row $r$ -- so the histogram-area algorithm run on
that exact $\textit{height}$ array is guaranteed to find the largest
rectangle with bottom edge on row $r$ specifically. Since every possible
all-$1$s rectangle has its bottom edge on \emph{some} row, running this
check for every row and taking the overall maximum is guaranteed not to
miss any candidate rectangle.

\subsection*{Algorithm}

\begin{enumerate}
  \item Initialize $\textit{height}[0..m-1] \gets 0$ and
        $\textit{maxArea} \gets 0$.
  \item For each row $r$ from $0$ to $n-1$:
  \begin{enumerate}
    \item For each column $j$: if $\textit{matrix}[r][j] = 1$:
          $\textit{height}[j] \gets \textit{height}[j]+1$; else:
          $\textit{height}[j] \gets 0$.
    \item Run the largest-rectangle-in-histogram algorithm on
          $\textit{height}$; update $\textit{maxArea}$ with the result.
  \end{enumerate}
  \item Return \textit{maxArea}.
\end{enumerate}

\subsection*{Dry run}

\dryrun{Take the matrix
$\begin{bmatrix} 1&0&1&0&0 \\ 1&0&1&1&1 \\ 1&1&1&1&1 \\ 1&0&0&1&0
\end{bmatrix}$.}

\begin{itemize}
  \item Row $0$: $\textit{height}=[1,0,1,0,0]$. Largest histogram
        rectangle: each isolated bar of height $1$, area $1$.
        $\textit{maxArea}=1$.
  \item Row $1$: $\textit{height}=[2,0,2,1,1]$. Largest rectangle: bars
        at indices $2,3,4$ with heights $2,1,1$ give area
        $1\times3=3$ (limited by the shortest, height $1$); also
        indices $0$ alone gives $2$; index $2$ alone gives $2$.
        Best here is $3$. $\textit{maxArea}=3$.
  \item Row $2$: $\textit{height}=[3,1,3,2,2]$. Largest rectangle:
        indices $2,3,4$ with heights $3,2,2$ give area $2\times3=6$.
        $\textit{maxArea}=6$.
  \item Row $3$: $\textit{height}=[4,0,0,3,0]$. Largest rectangle:
        index $0$ alone gives $4$; index $3$ alone gives $3$. Best here
        is $4$, not exceeding $6$.
\end{itemize}

The maximal rectangle has area $\boxed{6}$.

\subsection*{C++ implementation}

\begin{lstlisting}
#include <bits/stdc++.h>
using namespace std;

int largestRectangleArea(vector<int> heights) {
    heights.push_back(0);
    int n = heights.size();
    stack<int> st;
    long long maxArea = 0;

    for (int i = 0; i < n; i++) {
        while (!st.empty() && heights[i] < heights[st.top()]) {
            int j = st.top();
            st.pop();
            int width = st.empty() ? i : i - st.top() - 1;
            maxArea = max(maxArea, (long long)heights[j] * width);
        }
        st.push(i);
    }
    return maxArea;
}

int maximalRectangle(vector<vector<int>>& matrix) {
    int n = matrix.size(), m = matrix[0].size();
    vector<int> height(m, 0);
    int maxArea = 0;

    for (int r = 0; r < n; r++) {
        for (int j = 0; j < m; j++) {
            height[j] = (matrix[r][j] == 1) ? height[j] + 1 : 0;
        }
        maxArea = max(maxArea, largestRectangleArea(height));
    }
    return maxArea;
}

int main() {
    vector<vector<int>> matrix = {
        {1,0,1,0,0},
        {1,0,1,1,1},
        {1,1,1,1,1},
        {1,0,0,1,0}
    };
    cout << "Maximal rectangle area: " << maximalRectangle(matrix) << endl;
    return 0;
}
\end{lstlisting}

\begin{complexitybox}
\textbf{Time Complexity.} Updating the height array for each row costs
$O(m)$, and running the histogram algorithm costs $O(m)$, for a total of
$O(m)$ per row across $n$ rows, giving
\[
  O(nm).
\]
\textbf{Space Complexity.} $O(m)$ for the running height array and the
histogram algorithm's stack.
\end{complexitybox}

\begin{takeawaybox}
A 2D ``maximal region with some property'' problem often decomposes into
``run a 1D algorithm on a derived array, once per row (or column), with
the derived array updated incrementally as you move between rows.'' Once
you have Largest Rectangle in Histogram, Maximal Rectangle requires no new
stack logic at all -- only the row-by-row height bookkeeping is new.
\end{takeawaybox}

\section{Stock Span Problem}
\label{sec:sq-stock-span}

\noindent We now turn to the final group of problems in this chapter: a
handful of practical implementation and design problems that each combine
stacks or queues with a hash map or other auxiliary structure to solve a
real-world-flavored task efficiently.

\problemstatement{We are given the price of a stock on successive days,
arriving one at a time (a stream). For each new day's price, compute its
\textbf{span}: the number of consecutive days, ending with and including
today, for which today's price is greater than or equal to each of those
days' prices.}

\begin{patternbox}
\emph{``For each new arrival, look backward and count a contiguous run
based on a comparison''} is the \emph{previous} greater/smaller-element
family from Section~\ref{sec:sq-next-smaller}, but with a twist: we need
a \textbf{count} (the span length), not just the nearest qualifying
index. The key efficiency idea is that, instead of storing every
individual day on the stack, we can store \textbf{(price, span) pairs}
and merge spans together as we pop -- avoiding ever re-scanning days that
have already been absorbed into a previous day's span.
\end{patternbox}

\subsection*{Understanding the problem carefully}

The span of today's price is $1$ (counting today itself) plus the spans
of every immediately preceding day whose price is $\le$ today's price,
chained backward as far as this holds. If we stored only raw prices on
the stack, computing this span naively would require popping one day at a
time and could cost $O(n)$ per query in the worst case, but with no way
to skip over runs that were already resolved. Storing the \emph{span}
\emph{alongside} each price lets us absorb an entire previously-resolved
run in a single pop, rather than one day at a time.

\begin{ideabox}[Stack of (Price, Span) Pairs, Merging on Pop]
Maintain a stack of $(\textit{price}, \textit{span})$ pairs. For each new
price $p$: initialize $\textit{currentSpan} \gets 1$. While the stack is
not empty and its top's price $\le p$: pop the top, and add its
\textit{span} to $\textit{currentSpan}$ (this whole already-resolved run
is absorbed in one step, not re-examined day by day). Push
$(p, \textit{currentSpan})$, and report $\textit{currentSpan}$ as today's
span.
\end{ideabox}

\subsection*{Why absorbing a popped pair's entire span (not just its price) is correct}

If the top of the stack is a pair $(\textit{price}', \textit{span}')$
with $\textit{price}' \le p$, then \emph{every} day within that pair's
already-computed span also has a price $\le \textit{price}' \le p$ (that
is precisely what the pair's span certifies, by the inductive
construction of the algorithm) -- so all of those days extend today's
span too, and we can add the whole count $\textit{span}'$ in one step
rather than re-deriving it. This amortizes the total work: each day is
pushed exactly once and, across its entire lifetime, can only be merged
away (absorbed into a later pair) once, giving linear total work despite
some individual days having large spans.

\subsection*{Algorithm}

\begin{enumerate}
  \item Initialize an empty stack of (price, span) pairs.
  \item For each new price $p$:
  \begin{enumerate}
    \item $\textit{currentSpan} \gets 1$.
    \item While the stack is not empty and its top's price $\le p$: pop
          $(\textit{price}', \textit{span}')$;
          $\textit{currentSpan} \gets \textit{currentSpan} +
          \textit{span}'$.
    \item Push $(p, \textit{currentSpan})$.
    \item Report $\textit{currentSpan}$.
  \end{enumerate}
\end{enumerate}

\subsection*{Dry run}

\dryrun{Prices arrive in order: $100, 80, 60, 70, 60, 75, 85$.}

\begin{itemize}
  \item $100$: stack empty. span$=1$. Push $(100,1)$. Report $1$.
  \item $80$: top $(100,1)$, price $100>80$, no merge. span$=1$. Push
        $(80,1)$. Report $1$.
  \item $60$: top $(80,1)$, $80>60$, no merge. span$=1$. Push $(60,1)$.
        Report $1$.
  \item $70$: top $(60,1)$, $60\le70$: merge, span$=1+1=2$, pop. New top
        $(80,1)$, $80>70$, stop. Push $(70,2)$. Report $2$.
  \item $60$: top $(70,2)$, $70>60$, no merge. span$=1$. Push $(60,1)$.
        Report $1$.
  \item $75$: top $(60,1)$, $60\le75$: merge, span$=1+1=2$, pop. Top
        $(70,2)$, $70\le75$: merge, span$=2+2=4$, pop. Top $(80,1)$,
        $80>75$, stop. Push $(75,4)$. Report $4$.
  \item $85$: top $(75,4)$, $75\le85$: merge, span$=1+4=5$, pop. Top
        $(80,1)$, $80\le85$: merge, span$=5+1=6$, pop. Top $(100,1)$,
        $100>85$, stop. Push $(85,6)$. Report $6$.
\end{itemize}

Spans: $[1,1,1,2,1,4,6]$.

\subsection*{C++ implementation}

\begin{lstlisting}
#include <bits/stdc++.h>
using namespace std;

class StockSpanner {
    stack<pair<int,int>> st; // (price, span)

public:
    int next(int price) {
        int span = 1;
        while (!st.empty() && st.top().first <= price) {
            span += st.top().second;
            st.pop();
        }
        st.push({price, span});
        return span;
    }
};

int main() {
    StockSpanner spanner;
    for (int p : {100, 80, 60, 70, 60, 75, 85}) {
        cout << spanner.next(p) << " ";
    }
    cout << endl;
    return 0;
}
\end{lstlisting}

\begin{complexitybox}
\textbf{Time Complexity.} Each day is pushed once and merged away
(popped) at most once across its entire lifetime, giving
\[
  O(1) \text{ amortized per day}, \qquad O(n) \text{ total for } n \text{
  days}.
\]
\textbf{Space Complexity.} $O(n)$ for the stack in the worst case
(strictly decreasing prices, where every day stays on the stack).
\end{complexitybox}

\begin{takeawaybox}
When a monotonic-stack problem needs a \emph{count} of absorbed elements
rather than just the nearest boundary, store that count directly inside
each stack entry and merge counts on pop -- this avoids re-deriving the
count from scratch every time and keeps the whole algorithm at amortized
$O(1)$ per arrival, even for a problem (like this one) where the answer is
a span length, not a position.
\end{takeawaybox}

\section{The Celebrity Problem}
\label{sec:sq-celebrity-problem}

\problemstatement{There are $n$ people at a party. A \textbf{celebrity}
is a person whom \textbf{everyone else knows}, but who \textbf{knows no
one else}. We have access to a function $\textit{knows}(a,b)$ that
returns whether person $a$ knows person $b$. Find the celebrity if one
exists (there is at most one), using as few calls to $\textit{knows}$ as
possible -- ideally $O(n)$, not the $O(n^2)$ a full pairwise check would
need.}

\begin{patternbox}
\emph{``Repeatedly eliminate one of two candidates based on a pairwise
comparison, narrowing down to a single survivor''} is a stack-driven
elimination tournament: push every person as a candidate, then repeatedly
pop two candidates and use a single $\textit{knows}$ call to eliminate
(rather than confirm) one of them, since the celebrity's defining
properties make it possible to always rule out at least one of any two
candidates with a single check.
\end{patternbox}

\subsection*{Understanding the problem carefully}

Take any two candidates $a$ and $b$, and ask $\textit{knows}(a,b)$.

\begin{itemize}
  \item If $a$ knows $b$: then $a$ cannot be the celebrity (a celebrity
        knows \emph{no one}), so $a$ is eliminated.
  \item If $a$ does not know $b$: then $b$ cannot be the celebrity (the
        celebrity must be known by \emph{everyone}, but $a$ does not
        know $b$), so $b$ is eliminated.
\end{itemize}

Either way, a single call eliminates exactly one of the two candidates --
never both, and never neither -- so repeatedly applying this to pairs of
remaining candidates narrows the field down to a single survivor using
only $n-1$ calls total.

\begin{ideabox}[Stack-Driven Elimination, Then One Verification Pass]
Push every person $0, \dots, n-1$ onto a stack. While more than one
candidate remains: pop two, $a$ and $b$. If $\textit{knows}(a,b)$: push
$b$ back (eliminate $a$). Else: push $a$ back (eliminate $b$). After this
process, exactly one candidate remains -- but it is only a
\emph{candidate}, not yet confirmed (it is possible that \emph{no}
celebrity exists at all). Verify the survivor by checking that
\emph{everyone else} knows them and that they know \emph{no one else} --
an additional $O(n)$ scan.
\end{ideabox}

\subsection*{Why elimination narrows correctly, but still needs a final check}

The elimination process guarantees that \emph{if} a celebrity exists, they
are never eliminated: a true celebrity, compared against any other
candidate $x$, always satisfies $\textit{knows}(\textit{celebrity}, x) =
\text{false}$ (the celebrity knows no one) and
$\textit{knows}(x, \textit{celebrity}) = \text{true}$ (everyone knows the
celebrity) -- so whichever of the two checks the algorithm happens to make,
the celebrity is never the one eliminated. However, if \emph{no}
celebrity exists, the elimination process still produces \emph{some}
survivor (it always narrows to exactly one candidate, regardless of
whether that candidate has the celebrity property) -- which is why an
explicit final verification pass is required to confirm or reject the
survivor.

\subsection*{Algorithm}

\begin{enumerate}
  \item Push all $n$ people onto a stack.
  \item While the stack has more than one element:
  \begin{enumerate}
    \item Pop $a$, then pop $b$.
    \item If $\textit{knows}(a,b)$: push $b$. Else: push $a$.
  \end{enumerate}
  \item Let $\textit{candidate}$ be the single remaining person.
  \item Verify: for every other person $x$, check
        $\textit{knows}(x, \textit{candidate}) = \text{true}$ and
        $\textit{knows}(\textit{candidate}, x) = \text{false}$. If all
        checks pass, return $\textit{candidate}$; otherwise, return
        ``no celebrity exists.''
\end{enumerate}

\subsection*{Dry run}

\dryrun{Take $n=4$ with $\textit{knows}$ such that person $2$ is the true
celebrity: everyone knows $2$, and $2$ knows no one.}

Stack (top to bottom, after pushing $0,1,2,3$): top is $3$.

\begin{itemize}
  \item Pop $a=3$, $b=2$. Check $\textit{knows}(3,2)$: true (everyone
        knows $2$). Push $b=2$ (eliminate $3$). Stack $=[0,1,2]$.
  \item Pop $a=2$, $b=1$. Check $\textit{knows}(2,1)$: false ($2$ knows
        no one). Push $a=2$ (eliminate $1$). Stack $=[0,2]$.
  \item Pop $a=2$, $b=0$. Check $\textit{knows}(2,0)$: false. Push
        $a=2$ (eliminate $0$). Stack $=[2]$.
\end{itemize}

Candidate $=2$. Verification: every other person knows $2$ (true), and
$2$ knows none of them (true). Confirmed: person $2$ is the celebrity.

\subsection*{C++ implementation}

\begin{lstlisting}
#include <bits/stdc++.h>
using namespace std;

int findCelebrity(int n, function<bool(int,int)> knows) {
    stack<int> st;
    for (int i = 0; i < n; i++) st.push(i);

    while (st.size() > 1) {
        int a = st.top(); st.pop();
        int b = st.top(); st.pop();

        if (knows(a, b)) st.push(b); // eliminate a
        else st.push(a);             // eliminate b
    }

    int candidate = st.top();
    for (int x = 0; x < n; x++) {
        if (x == candidate) continue;
        if (!knows(x, candidate) || knows(candidate, x)) {
            return -1; // no celebrity
        }
    }
    return candidate;
}

int main() {
    // Person 2 is the celebrity: everyone knows 2, 2 knows no one.
    vector<vector<bool>> M = {
        {0,1,1,1},
        {0,0,1,0},
        {0,0,0,0},
        {0,1,1,0}
    };
    auto knows = [&](int a, int b) { return M[a][b]; };

    cout << "Celebrity: " << findCelebrity(4, knows) << endl;
    return 0;
}
\end{lstlisting}

\begin{complexitybox}
\textbf{Time Complexity.} The elimination phase makes exactly $n-1$
calls to $\textit{knows}$ (one per pop-pair, each removing one candidate
from a starting pool of $n$), and the verification phase makes up to
$2(n-1)$ more calls, giving an overall time complexity of
\[
  O(n).
\]
\textbf{Space Complexity.} $O(n)$ for the stack.
\end{complexitybox}

\begin{takeawaybox}
When a problem guarantees that comparing any two candidates always lets
you safely eliminate \emph{at least one} of them (never both, never
neither), a stack-driven elimination tournament narrows $n$ candidates
to $1$ in $O(n)$ comparisons. This is a different flavor of stack usage
from the monotonic-stack family -- here the stack just manages an
elimination order, not a sorted skyline -- and it is always worth pairing
with an explicit final verification step whenever the narrowing process
does not, by itself, guarantee the survivor actually has the desired
property.
\end{takeawaybox}

\section{Sliding Window Maximum}
\label{sec:sq-sliding-window-max}

\problemstatement{Given an array $\textit{nums}$ and a window size $k$,
return an array containing the maximum element of every contiguous
window of size $k$ as it slides from the start of the array to the end.}

\begin{patternbox}
\emph{``Maximum (or minimum) over every sliding window''} is the
\textbf{monotonic deque} keyword from Section~\ref{sec:sq-intro}: a
single index can remain a useful candidate for the maximum across
\emph{many} consecutive windows, but it must also eventually be evicted
once the window slides past it -- requiring removal from \textbf{both}
ends (the back, when a better candidate arrives; the front, when the
oldest index falls out of the window), which is exactly what a
double-ended queue (deque) provides.
\end{patternbox}

\subsection*{Understanding the problem carefully}

Maintain a deque of \emph{indices} (not values), kept in decreasing order
of their corresponding values from front to back, and restricted to
indices currently inside the window. The front of the deque is always the
index of the current window's maximum, because it is both the largest
remaining value (by the decreasing-order invariant) and guaranteed to
still be within the window (indices that fall out are actively removed
from the front).

\begin{ideabox}[Monotonic Deque: Evict from the Back on Value, from the Front on Position]
For each new index $i$: first, while the deque is non-empty and the value
at its back index is $\le \textit{nums}[i]$, pop from the back (those
indices can never again be the maximum of any window containing $i$,
since $i$ is both newer and at least as large -- the same domination
argument as next greater element). Push $i$ to the back. Then, while the
deque's front index is $\le i-k$ (i.e.\ has fallen outside the current
window), pop from the front. Once $i \ge k-1$, the deque's front index is
the maximum of the current window.
\end{ideabox}

\subsection*{Why eviction from both ends preserves correctness}

Eviction from the \textbf{back} is exactly the monotonic-stack domination
argument from Section~\ref{sec:sq-next-greater-1}: a smaller, older index
can never be the maximum once a newer, at-least-as-large value has
arrived, anywhere within the window. Eviction from the \textbf{front} is
new to this section and is required because, unlike a plain monotonic
stack, indices here have a finite lifetime determined by the window's
position -- even an index that has never been dominated by value must
still be discarded once it ages out of the window. Both evictions
together guarantee that the deque, at every step, contains exactly the
indices that are both currently in-window \emph{and} not dominated by
any later in-window value, kept in decreasing order -- so its front is
always correct.

\subsection*{Algorithm}

\begin{enumerate}
  \item Initialize an empty deque of indices and an empty result list.
  \item For $i$ from $0$ to $n-1$:
  \begin{enumerate}
    \item While the deque is non-empty and
          $\textit{nums}[\text{back index}] \le \textit{nums}[i]$: pop
          from the back.
    \item Push $i$ to the back.
    \item While the deque's front index $\le i-k$: pop from the front.
    \item If $i \ge k-1$: append $\textit{nums}[\text{front index}]$ to
          the result.
  \end{enumerate}
  \item Return the result.
\end{enumerate}

\subsection*{Dry run}

\dryrun{Take $\textit{nums} = [1,3,-1,-3,5,3,6,7]$, $k=3$.}

\begin{itemize}
  \item $i=0$ ($1$): deque empty, push. Deque $=[0]$. $i<k-1=2$, no
        output.
  \item $i=1$ ($3$): back is index $0$ (value $1\le3$), pop. Push $1$.
        Deque $=[1]$. $i<2$, no output.
  \item $i=2$ ($-1$): back is index $1$ (value $3 \not\le -1$), no pop.
        Push $2$. Deque $=[1,2]$. Front index $1 > i-k=-1$, no pop.
        $i=2\ge2$: output $\textit{nums}[1]=3$.
  \item $i=3$ ($-3$): no pop (back value $-1\not\le-3$). Push $3$.
        Deque $=[1,2,3]$. Front index $1 \le i-k=0$? No ($1>0$). Output
        $\textit{nums}[1]=3$.
  \item $i=4$ ($5$): back values $-3,-1,3$ all $\le5$: pop all. Deque
        $=[]$, push $4$. Deque $=[4]$. Front $4 \le i-k=1$? No. Output
        $\textit{nums}[4]=5$.
  \item $i=5$ ($3$): back value $5\not\le3$, no pop. Push $5$. Deque
        $=[4,5]$. Front $4\le i-k=2$? No. Output $\textit{nums}[4]=5$.
  \item $i=6$ ($6$): back values $3,5$ both $\le6$: pop both. Push $6$.
        Deque $=[6]$. Front $6\le i-k=3$? No. Output
        $\textit{nums}[6]=6$.
  \item $i=7$ ($7$): back value $6\le7$, pop. Push $7$. Deque $=[7]$.
        Front $7\le i-k=4$? No. Output $\textit{nums}[7]=7$.
\end{itemize}

Result: $[3,3,5,5,6,7]$.

\subsection*{C++ implementation}

\begin{lstlisting}
#include <bits/stdc++.h>
using namespace std;

vector<int> maxSlidingWindow(vector<int>& nums, int k) {
    deque<int> dq; // stores indices
    vector<int> result;

    for (int i = 0; i < (int)nums.size(); i++) {
        while (!dq.empty() && nums[dq.back()] <= nums[i]) {
            dq.pop_back();
        }
        dq.push_back(i);

        while (dq.front() <= i - k) {
            dq.pop_front();
        }

        if (i >= k - 1) {
            result.push_back(nums[dq.front()]);
        }
    }
    return result;
}

int main() {
    vector<int> nums = {1, 3, -1, -3, 5, 3, 6, 7};
    for (int x : maxSlidingWindow(nums, 3)) cout << x << " ";
    cout << endl;
    return 0;
}
\end{lstlisting}

\begin{complexitybox}
\textbf{Time Complexity.} Each index is pushed to the back once and
popped (from either end) at most once across the whole scan, giving
\[
  O(n).
\]
\textbf{Space Complexity.} $O(k)$ for the deque in the worst case.
\end{complexitybox}

\begin{takeawaybox}
A monotonic deque is a monotonic stack with one extra capability: aging
out elements from the opposite end once they fall outside a window. Any
time a problem combines ``sliding window'' with ``maximum/minimum within
the window,'' this is the data structure to reach for, rather than
re-scanning each window from scratch (which would cost $O(nk)$) or using
a heap (which would cost $O(n \log k)$ but require extra bookkeeping to
evict stale entries lazily).
\end{takeawaybox}

\section{LRU Cache}
\label{sec:sq-lru-cache}

\problemstatement{Design a cache with a fixed capacity supporting
$\textit{get}(\textit{key})$ (return the value if present, else $-1$, and
mark the key as recently used) and $\textit{put}(\textit{key},
\textit{value})$ (insert or update a key, marking it recently used; if the
cache is full and a new key must be inserted, evict the
\textbf{least recently used} key first). Both operations must run in
$O(1)$.}

\begin{patternbox}
\emph{``Most/least recently used, with fast lookup''} is the keyword
combination from Section~\ref{sec:sq-intro}: $O(1)$ lookup by key demands
a hash map, while $O(1)$ ``move this to the most-recent end'' and
``evict the least-recent end'' demands a structure with fast removal and
insertion at \emph{both} ends and from the \emph{middle} -- which a plain
array or singly linked list cannot offer, but a \textbf{doubly linked
list} can.
\end{patternbox}

\subsection*{Understanding the problem carefully}

A hash map alone gives $O(1)$ lookup but no notion of ``recency
order.'' A queue alone gives recency order but $O(n)$ lookup by key, and
moving an arbitrary key to the front of recency would require knowing its
position. Combining both is the standard fix: a \textbf{doubly linked
list} maintains recency order (most recently used at one end, least
recently used at the other), and a \textbf{hash map from key to that
key's node} gives $O(1)$ lookup -- crucially, because the list is
\emph{doubly} linked, removing a node found via the hash map (to move it
or evict it) takes $O(1)$, since we have direct pointers to its neighbors
without needing to traverse from either end.

\begin{ideabox}[Doubly Linked List + Hash Map]
Maintain a doubly linked list with two sentinel nodes, \textit{head} and
\textit{tail}, where the side nearest \textit{head} represents
\emph{most} recently used and the side nearest \textit{tail} represents
\emph{least} recently used. A hash map stores $\textit{key} \to
\textit{node pointer}$. On $\textit{get}(\textit{key})$: if present,
unlink the node from its current position and re-insert it right after
\textit{head} (marking it most recent), then return its value; if absent,
return $-1$. On $\textit{put}(\textit{key}, \textit{value})$: if the key
exists, update its value and move its node to just after \textit{head}
(same as a successful get). If the key is new and the cache is at
capacity, remove the node just before \textit{tail} (the least recently
used) along with its hash map entry; then insert a new node for this key
just after \textit{head}.
\end{ideabox}

\subsection*{Why doubly linked, not singly linked, is essential}

The operation that makes this problem hard is not insertion or removal at
the \emph{ends} (a singly linked list with head and tail pointers could
handle that, as in Section~\ref{sec:sq-queue-linked-list}) -- it is
removal from an \textbf{arbitrary middle position}, which happens every
time an already-cached key is accessed again and must be promoted to
most-recently-used. Removing a node from the middle of a singly linked
list requires knowing its \emph{predecessor}, which means traversing from
the head -- an $O(n)$ operation. A doubly linked list stores each node's
predecessor pointer directly, so the hash map's stored node pointer alone
is enough to unlink it in $O(1)$, with no traversal needed at all.

\subsection*{Algorithm}

\begin{enumerate}
  \item Maintain sentinel \textit{head} and \textit{tail} nodes, a hash
        map $\textit{key} \to \textit{node}$, and a capacity limit.
  \item \textsc{Get}($\textit{key}$): if absent, return $-1$. Else:
        unlink the node; re-insert it right after \textit{head}; return
        its value.
  \item \textsc{Put}($\textit{key}, \textit{value}$): if present, update
        the node's value and move it to right after \textit{head}
        (as in \textsc{Get}). If absent and at capacity: remove the
        node just before \textit{tail} (and its map entry). Create a new
        node, insert it right after \textit{head}, and add it to the map.
\end{enumerate}

\subsection*{Dry run}

\dryrun{Capacity $2$. \textsc{Put}(1,1), \textsc{Put}(2,2),
\textsc{Get}(1), \textsc{Put}(3,3) (evicts LRU), \textsc{Get}(2).}

\begin{itemize}
  \item Put(1,1): list (MRU$\to$LRU) $=[1]$.
  \item Put(2,2): list $=[2,1]$.
  \item Get(1): found; move to front. List $=[1,2]$. Returns $1$.
  \item Put(3,3): cache full ($2$ keys); evict LRU end ($2$). List
        $=[3,1]$.
  \item Get(2): key $2$ was evicted; returns $-1$.
\end{itemize}

\subsection*{C++ implementation}

\begin{lstlisting}
#include <bits/stdc++.h>
using namespace std;

class LRUCache {
    struct Node {
        int key, value;
        Node* prev = nullptr;
        Node* next = nullptr;
        Node(int k, int v) : key(k), value(v) {}
    };

    int capacity;
    Node* head; // sentinel: head->next is most recently used
    Node* tail; // sentinel: tail->prev is least recently used
    unordered_map<int, Node*> map;

    void remove(Node* node) {
        node->prev->next = node->next;
        node->next->prev = node->prev;
    }

    void insertAfterHead(Node* node) {
        node->next = head->next;
        node->prev = head;
        head->next->prev = node;
        head->next = node;
    }

public:
    LRUCache(int cap) : capacity(cap) {
        head = new Node(-1, -1);
        tail = new Node(-1, -1);
        head->next = tail;
        tail->prev = head;
    }

    int get(int key) {
        if (!map.count(key)) return -1;
        Node* node = map[key];
        remove(node);
        insertAfterHead(node);
        return node->value;
    }

    void put(int key, int value) {
        if (map.count(key)) {
            Node* node = map[key];
            node->value = value;
            remove(node);
            insertAfterHead(node);
            return;
        }
        if ((int)map.size() == capacity) {
            Node* lru = tail->prev;
            remove(lru);
            map.erase(lru->key);
            delete lru;
        }
        Node* node = new Node(key, value);
        insertAfterHead(node);
        map[key] = node;
    }
};

int main() {
    LRUCache cache(2);
    cache.put(1, 1);
    cache.put(2, 2);
    cout << cache.get(1) << endl; // 1
    cache.put(3, 3); // evicts key 2
    cout << cache.get(2) << endl; // -1
    return 0;
}
\end{lstlisting}

\begin{complexitybox}
\textbf{Time Complexity.} $O(1)$ for both $\textit{get}$ and
$\textit{put}$ -- hash map lookup, plus a constant number of pointer
updates for the doubly linked list operations.
\textbf{Space Complexity.} $O(\textit{capacity})$ for the nodes and hash
map.
\end{complexitybox}

\begin{takeawaybox}
``Fast lookup by key'' plus ``fast reordering by recency, including
removal from the middle'' is the signature combination that demands a
hash map paired with a \emph{doubly} linked list, not a singly linked one
-- the doubly linked structure is what converts an $O(n)$ traversal-to-find-predecessor
into an $O(1)$ direct unlink. The same pairing, with a different
eviction rule, solves LFU Cache next.
\end{takeawaybox}

\section{LFU Cache}
\label{sec:sq-lfu-cache}

\problemstatement{Design a cache with a fixed capacity supporting
$\textit{get}(\textit{key})$ and $\textit{put}(\textit{key},
\textit{value})$, both in $O(1)$, where eviction (when the cache is full
and a new key must be inserted) removes the \textbf{least frequently
used} key; if multiple keys share the lowest frequency, the
\textbf{least recently used} among them is evicted.}

\begin{patternbox}
LFU is LRU (Section~\ref{sec:sq-lru-cache}) with a second tie-breaking
dimension added: eviction priority is now \emph{(frequency, recency)},
not recency alone. The keyword to notice: whenever a tie-breaking rule is
layered on top of a structure you already have working (here, LRU's
hash-map-plus-doubly-linked-list), try \textbf{bucketing} by the primary
key (frequency) and reusing the existing structure (an LRU list) as the
secondary tie-breaker \emph{within} each bucket, rather than designing an
entirely new structure from scratch.
\end{patternbox}

\subsection*{Understanding the problem carefully}

Group keys by their current access frequency into separate buckets. Each
bucket, internally, only ever needs to break ties by recency among keys
that share that exact frequency -- which is exactly the LRU Cache
substructure from Section~\ref{sec:sq-lru-cache}, reused once per
distinct frequency value. Eviction always targets the \emph{lowest}
currently-occupied frequency bucket, and within it, removes that
bucket's least-recently-used key (the LRU substructure's own eviction
rule, applied locally).

\begin{ideabox}[A Hash Map of Frequency Buckets, Each an LRU List]
Maintain: a hash map $\textit{keyData}$ from key to $(\textit{value},
\textit{freq})$; a hash map $\textit{freqBuckets}$ from frequency to a
doubly linked list (an LRU list) of keys currently at that frequency,
ordered by recency within the bucket; and a single integer
$\textit{minFreq}$, tracking the smallest frequency currently occupied by
any key. On every access to an existing key (via \textit{get} or an
updating \textit{put}): remove it from its current frequency bucket
(inserting it at the front, i.e.\ most-recent position, of the
\emph{next} frequency bucket up), incrementing its stored frequency; if
its old bucket becomes empty and was the minimum, increment
$\textit{minFreq}$. On inserting a brand-new key when the cache is full:
evict the least-recently-used key from the $\textit{minFreq}$ bucket, then
insert the new key into the frequency-$1$ bucket and reset
$\textit{minFreq} \gets 1$.
\end{ideabox}

\subsection*{Why tracking \textit{minFreq} explicitly avoids an expensive search}

Without an explicit running $\textit{minFreq}$, finding ``the lowest
currently-occupied frequency'' would require scanning all buckets,
costing more than $O(1)$. Because frequency only ever \textbf{increases}
by exactly $1$ per access, and a brand-new key always enters at frequency
$1$, we can maintain $\textit{minFreq}$ incrementally with simple,
local updates: it can only decrease (to $1$) when a new key is inserted,
and it can only increase when the bucket it currently points to becomes
completely empty (because the key that just left was the last one at that
frequency, and every other key's frequency is, by definition, at least
\textit{minFreq}, so the new minimum must be \textit{minFreq}+1 in that
specific case). This local bookkeeping replaces a potentially expensive
global search with $O(1)$ updates.

\subsection*{Algorithm}

\begin{enumerate}
  \item \textsc{Touch}($\textit{key}$) (shared helper for get and
        updating put): look up $(\textit{value}, f)$ in
        $\textit{keyData}$; remove $\textit{key}$ from
        $\textit{freqBuckets}[f]$; if that bucket is now empty and
        $f = \textit{minFreq}$, increment \textit{minFreq}; insert
        $\textit{key}$ at the front of $\textit{freqBuckets}[f+1]$;
        update $\textit{keyData}[\textit{key}] \gets (\textit{value},
        f+1)$.
  \item \textsc{Get}($\textit{key}$): if absent, return $-1$. Else:
        \textsc{Touch}($\textit{key}$); return its value.
  \item \textsc{Put}($\textit{key}, \textit{value}$): if present, update
        the stored value and call \textsc{Touch}($\textit{key}$).
        Otherwise, if at capacity: evict the least-recently-used key
        from $\textit{freqBuckets}[\textit{minFreq}]$ (and its
        $\textit{keyData}$ entry). Insert $\textit{key}$ at the front of
        $\textit{freqBuckets}[1]$; set $\textit{keyData}[\textit{key}]
        \gets (\textit{value}, 1)$; set $\textit{minFreq} \gets 1$.
\end{enumerate}

\subsection*{Dry run}

\dryrun{Capacity $2$. \textsc{Put}(1,1), \textsc{Put}(2,2),
\textsc{Get}(1) (freq of $1$ becomes $2$), \textsc{Put}(3,3) (cache full;
evict lowest-frequency key, which is $2$, frequency $1$), \textsc{Get}(2).}

\begin{itemize}
  \item Put(1,1): bucket[1]$=[1]$. minFreq$=1$.
  \item Put(2,2): bucket[1]$=[2,1]$ (2 most recent). minFreq$=1$.
  \item Get(1): touch $1$ (freq $1\to2$). Remove from bucket[1]:
        bucket[1]$=[2]$ (not empty, minFreq stays $1$). Insert into
        bucket[2]$=[1]$.
  \item Put(3,3): cache full ($2$ keys). Evict from bucket[\textit{minFreq}=1]'s
        LRU end: bucket[1]$=[2]$, evict $2$. bucket[1] now empty. Insert
        $3$ into bucket[1]$=[3]$. minFreq reset to $1$.
  \item Get(2): key $2$ was evicted; returns $-1$.
\end{itemize}

\subsection*{C++ implementation}

\begin{lstlisting}
#include <bits/stdc++.h>
using namespace std;

class LFUCache {
    int capacity, minFreq;
    unordered_map<int, pair<int,int>> keyData; // key -> (value, freq)
    unordered_map<int, list<int>> freqBuckets;  // freq -> keys (front = most recent)
    unordered_map<int, list<int>::iterator> keyIter; // key -> its position in its bucket

    void touch(int key) {
        auto [value, freq] = keyData[key];
        freqBuckets[freq].erase(keyIter[key]);
        if (freqBuckets[freq].empty()) {
            freqBuckets.erase(freq);
            if (freq == minFreq) minFreq++;
        }
        freqBuckets[freq + 1].push_front(key);
        keyIter[key] = freqBuckets[freq + 1].begin();
        keyData[key] = {value, freq + 1};
    }

public:
    LFUCache(int cap) : capacity(cap), minFreq(0) {}

    int get(int key) {
        if (!keyData.count(key)) return -1;
        touch(key);
        return keyData[key].first;
    }

    void put(int key, int value) {
        if (capacity == 0) return;

        if (keyData.count(key)) {
            keyData[key].first = value;
            touch(key);
            return;
        }

        if ((int)keyData.size() == capacity) {
            int evictKey = freqBuckets[minFreq].back();
            freqBuckets[minFreq].pop_back();
            keyData.erase(evictKey);
            keyIter.erase(evictKey);
        }

        keyData[key] = {value, 1};
        freqBuckets[1].push_front(key);
        keyIter[key] = freqBuckets[1].begin();
        minFreq = 1;
    }
};

int main() {
    LFUCache cache(2);
    cache.put(1, 1);
    cache.put(2, 2);
    cout << cache.get(1) << endl; // 1
    cache.put(3, 3);              // evicts key 2
    cout << cache.get(2) << endl; // -1
    return 0;
}
\end{lstlisting}

\begin{complexitybox}
\textbf{Time Complexity.} $O(1)$ for both $\textit{get}$ and
$\textit{put}$ -- every step is a hash map lookup or a constant number of
doubly linked list operations, with $\textit{minFreq}$ maintained
incrementally rather than recomputed.
\textbf{Space Complexity.} $O(\textit{capacity})$ across the hash maps
and linked lists.
\end{complexitybox}

\begin{takeawaybox}
This closing problem of the chapter is a deliberate capstone, much like
Median of a Row-wise Sorted Matrix was for the Binary Search chapter: it
requires recognizing that a new, seemingly more complex requirement
(frequency-based eviction) can be built by \emph{bucketing} on the new
dimension and reusing an already-solved structure (LRU's doubly linked
list) as the tie-breaker within each bucket -- composition of known
pieces, rather than a wholly new design, solves the harder problem.
\end{takeawaybox}

